% Môn: Vật lý
% Lớp: 12
% Chương: Chương II. Khí lí tưởng
% Bài: L12C2 Bài 10. Định luật Charles
% Số câu: 170
% App đọc file này trực tiếp — sửa rồi Commit trên GitHub, bấm Đồng bộ GitHub .tex

% L12C2 Bài 10. Định luật Charles

% ===== Câu 1 =====
% ID: L12C2B10-01-DS
% Mức: TH
\begin{ex}
% ID: L12C2B10-01-DS
% Mức: TH
Một bình thép kín có dung tích $10\ \mathrm{l}$ chứa khí hydrogen ở áp suất ban đầu là $2 \cdot 10^5\ \mathrm{Pa}$. Người ta dùng một bình tích năng lượng để nén thêm khí hydrogen vào bình thép. Sau một thời gian, áp suất tuyệt đối trong bình thép tăng lên đến $5 \cdot 10^5\ \mathrm{Pa}$. Biết quá trình diễn ra đẳng nhiệt.
\choiceTF
{Thể tích khí hydrogen lúc sau trong bình thép tăng lên gấp $2,5\ \mathrm{lần}$ vì áp suất tăng.}
{\True Vì bình bằng thép kín nên thể tích của lượng khí chứa trong bình lúc sau vẫn bằng $10\ \mathrm{l}$.}
{\True Thể tích của lượng khí nén thêm vào bình thép nếu quy về áp suất khí quyển $10^5\ \mathrm{Pa}$ có giá trị bằng $30\ \mathrm{l}$.}
{Khối lượng khí chứa trong bình thép hoàn toàn không thay đổi vì đây là bình thép kín.}
\loigiai{
\begin{itemchoice}
\item (Sai) Thể tích khí hydrogen lúc sau trong bình thép tăng lên gấp $2,5\ \mathrm{lần}$ vì áp suất tăng. (Vì): Vì bình bằng thép kín có thể tích không đổi nên thể tích khí luôn bằng dung tích bình là $10\ \mathrm{l}$ chứ không thể tăng lên. Mệnh đề sai
\item (Đúng) Vì bình bằng thép kín nên thể tích của lượng khí chứa trong bình lúc sau vẫn bằng $10\ \mathrm{l}$. (Vì): Bình thép có dung tích cố định nên thể tích khí sau khi nén thêm vẫn bằng $10\ \mathrm{l}$. Mệnh đề đúng
\item (Đúng) Thể tích của lượng khí nén thêm vào bình thép nếu quy về áp suất khí quyển $10^5\ \mathrm{Pa}$ có giá trị bằng $30\ \mathrm{l}$. (Vì): Áp dụng định luật Boyle cho quá trình tích áp, thể tích khí thêm vào quy về áp suất $p_0 = 10^5\ \mathrm{Pa}$ là:. Mệnh đề đúng
\item (Sai) Khối lượng khí chứa trong bình thép hoàn toàn không thay đổi vì đây là bình thép kín. (Vì): Vì người ta nén thêm một lượng khí hydrogen từ bên ngoài vào bình nên khối lượng khí trong bình phải tăng lên chứ không giữ nguyên. Mệnh đề sai
\end{itemchoice}
}
\end{ex}

% ===== Câu 2 =====
% ID: L12C2B10-01-TL
% Mức: TH
\dangbt{Bài toán tổng hợp về quá trình đẳng áp}
\begin{ex}
% ID: L12C2B10-01-TL
% Mức: TH
Trong dạng “Bài toán tổng hợp về quá trình đẳng áp”, Trình bày cơ sở lí thuyết của dạng “Bài toán tổng hợp về quá trình đẳng áp” và nêu điều kiện áp dụng.
\choiceTF

\loigiai{
Nêu đúng đại lượng không đổi, hệ thức đặc trưng và yêu cầu dùng nhiệt độ kelvin, áp suất tuyệt đối khi có liên quan.
}
\end{ex}

% ===== Câu 3 =====
% ID: L12C2B10-01-TLN
% Mức: TH
\begin{ex}
% ID: L12C2B10-01-TLN
% Mức: TH
Trong dạng “Bài toán tổng hợp về quá trình đẳng áp”, Khí 3 lít ở 300 $K$ được đun đẳng áp đến 360 K. Tính thể tích sau theo lít.
\shortans{3,6}
\loigiai{
$V_2=3\cdot360/300=3,6$ lít.
}
\end{ex}

% ===== Câu 4 =====
% ID: L12C2B10-01-TN
% Mức: NB
\begin{ex}
% ID: L12C2B10-01-TN
% Mức: NB
Trong dạng “Bài toán tổng hợp về quá trình đẳng áp”, Trong quá trình đẳng áp của một lượng khí xác định, đại lượng không đổi là
\choice
{\True áp suất}
{thể tích}
{nhiệt độ}
{khối lượng riêng}
\loigiai{
Đáp án A. Quá trình đẳng áp có áp suất không đổi.
}
\end{ex}

% ===== Câu 5 =====
% ID: L12C2B10-02-DS
% Mức: TH
\begin{ex}
% ID: L12C2B10-02-DS
% Mức: TH
Một bình thép kín có dung tích $10\ \mathrm{l}$ chứa khí hydrogen ở áp suất ban đầu là $2 \cdot 10^5\ \mathrm{Pa}$. Người ta dùng một bình tích năng lượng để nén thêm khí hydrogen vào bình thép. Sau một thời gian, áp suất tuyệt đối trong bình thép tăng lên đến $5 \cdot 10^5\ \mathrm{Pa}$. Biết quá trình diễn ra đẳng nhiệt.
\choiceTF
{Thể tích khí hydrogen lúc sau trong bình thép tăng lên gấp $2,5\ \mathrm{lần}$ vì áp suất tăng.}
{\True Vì bình bằng thép kín nên thể tích của lượng khí chứa trong bình lúc sau vẫn bằng $10\ \mathrm{l}$.}
{\True Thể tích của lượng khí nén thêm vào bình thép nếu quy về áp suất khí quyển $10^5\ \mathrm{Pa}$ có giá trị bằng $30\ \mathrm{l}$.}
{Khối lượng khí chứa trong bình thép hoàn toàn không thay đổi vì đây là bình thép kín.}
\loigiai{
\begin{itemchoice}
\item (Sai) Thể tích khí hydrogen lúc sau trong bình thép tăng lên gấp $2,5\ \mathrm{lần}$ vì áp suất tăng. (Vì): Vì bình bằng thép kín có thể tích không đổi nên thể tích khí luôn bằng dung tích bình là $10\ \mathrm{l}$ chứ không thể tăng lên. Mệnh đề sai
\item (Đúng) Vì bình bằng thép kín nên thể tích của lượng khí chứa trong bình lúc sau vẫn bằng $10\ \mathrm{l}$. (Vì): Bình thép có dung tích cố định nên thể tích khí sau khi nén thêm vẫn bằng $10\ \mathrm{l}$. Mệnh đề đúng
\item (Đúng) Thể tích của lượng khí nén thêm vào bình thép nếu quy về áp suất khí quyển $10^5\ \mathrm{Pa}$ có giá trị bằng $30\ \mathrm{l}$. (Vì): Áp dụng định luật Boyle cho quá trình tích áp, thể tích khí thêm vào quy về áp suất $p_0 = 10^5\ \mathrm{Pa}$ là:. Mệnh đề đúng
\item (Sai) Khối lượng khí chứa trong bình thép hoàn toàn không thay đổi vì đây là bình thép kín. (Vì): Vì người ta nén thêm một lượng khí hydrogen từ bên ngoài vào bình nên khối lượng khí trong bình phải tăng lên chứ không giữ nguyên. Mệnh đề sai
\end{itemchoice}
}
\end{ex}

% ===== Câu 6 =====
% ID: L12C2B10-02-TL
% Mức: TH
\begin{ex}
% ID: L12C2B10-02-TL
% Mức: TH
Trong dạng “Bài toán tổng hợp về quá trình đẳng áp”, Mô tả dạng đồ thị đặc trưng và cách nhận biết quá trình trong dạng “Bài toán tổng hợp về quá trình đẳng áp”.
\choiceTF

\loigiai{
Xác định các trục, đại lượng không đổi và suy ra hình dạng đồ thị từ hệ thức vật lí tương ứng.
}
\end{ex}

% ===== Câu 7 =====
% ID: L12C2B10-02-TLN
% Mức: TH
\begin{ex}
% ID: L12C2B10-02-TLN
% Mức: TH
Trong dạng “Bài toán tổng hợp về quá trình đẳng áp”, Khí có thể tích 5 lít ở 310 $K$, đẳng áp tăng thể tích lên 15 lít. Tính nhiệt độ sau theo K.
\shortans{930}
\loigiai{
$T_2=T_1V_2/V_1=310\cdot3=930$ K.
}
\end{ex}

% ===== Câu 8 =====
% ID: L12C2B10-02-TN
% Mức: NB
\begin{ex}
% ID: L12C2B10-02-TN
% Mức: NB
Trong dạng “Bài toán tổng hợp về quá trình đẳng áp”, Hệ thức đúng của định luật Charles là
\choice
{\True $V_1/T_1=V_2/T_2$}
{$p_1V_1=p_2V_2$}
{$p_1/T_1=p_2/T_2$}
{$V_1T_1=V_2T_2$}
\loigiai{
Đáp án A. Ở áp suất không đổi, $V/T$ là hằng số.
}
\end{ex}

% ===== Câu 9 =====
% ID: L12C2B10-03-DS
% Mức: TH
\begin{ex}
% ID: L12C2B10-03-DS
% Mức: TH
Xét các phát biểu sau trong dạng “Bài toán tổng hợp về quá trình đẳng áp”.
\choiceTF
{\True Trong quá trình đẳng áp của một lượng khí xác định, đại lượng không đổi là — phát biểu đúng là: áp suất.}
{Trong quá trình đẳng áp của một lượng khí xác định, đại lượng không đổi là — phát biểu thể tích là đúng.}
{\True Hệ thức đúng của định luật Charles là — phát biểu đúng là: $V_1/T_1=V_2/T_2$.}
{Hệ thức đúng của định luật Charles là — phát biểu $p_1/T_1=p_2/T_2$ là đúng.}
\loigiai{
\begin{itemchoice}
\item (Đúng) Trong quá trình đẳng áp của một lượng khí xác định, đại lượng không đổi là — phát biểu đúng là: áp suất. (Vì): Quá trình đẳng áp có áp suất không đổi. b) Sai: trái với hệ thức hoặc điều kiện đã nêu. c) Đúng: Ở áp suất không đổi, $V/T$ là hằng số. d) Sai: trái với hệ thức hoặc điều kiện đã nêu
\item (Sai) Trong quá trình đẳng áp của một lượng khí xác định, đại lượng không đổi là — phát biểu thể tích là đúng.
\item (Đúng) Hệ thức đúng của định luật Charles là — phát biểu đúng là: $V_1/T_1=V_2/T_2$.
\item (Sai) Hệ thức đúng của định luật Charles là — phát biểu $p_1/T_1=p_2/T_2$ là đúng.
\end{itemchoice}
}
\end{ex}

% ===== Câu 10 =====
% ID: L12C2B10-03-TL
% Mức: VD
\begin{bt}
% ID: L12C2B10-03-TL
% Mức: VD
Trong dạng “Bài toán tổng hợp về quá trình đẳng áp”, Nhiệt độ tuyệt đối tăng từ 320 $K$ lên 380 $K$ ở áp suất không đổi. Tính $V_2/V_1$.
\loigiai{
$V_2/V_1=T_2/T_1=1,19$.
}
\end{bt}

% ===== Câu 11 =====
% ID: L12C2B10-03-TLN
% Mức: VD
\begin{ex}
% ID: L12C2B10-03-TLN
% Mức: VD
Trong dạng “Bài toán tổng hợp về quá trình đẳng áp”, Nhiệt độ tuyệt đối tăng từ 320 $K$ lên 380 $K$ ở áp suất không đổi. Tính $V_2/V_1$.
\shortans{1,19}
\loigiai{
$V_2/V_1=T_2/T_1=1,19$.
}
\end{ex}

% ===== Câu 12 =====
% ID: L12C2B10-03-TN
% Mức: TH
\begin{ex}
% ID: L12C2B10-03-TN
% Mức: TH
Trong dạng “Bài toán tổng hợp về quá trình đẳng áp”, Khí có thể tích 4 lít ở 310 $K$, đun đẳng áp đến 380 K. Thể tích sau gần bằng
\choice
{\True 4,9 lít}
{3,26 lít}
{4 lít}
{9,81 lít}
\loigiai{
Đáp án A. $V_2=V_1T_2/T_1=4\cdot380/310=4,9$ lít.
}
\end{ex}

% ===== Câu 13 =====
% ID: L12C2B10-04-DS
% Mức: TH
\begin{ex}
% ID: L12C2B10-04-DS
% Mức: TH
Xét các phát biểu sau trong dạng “Bài toán tổng hợp về quá trình đẳng áp”.
\choiceTF
{\True Khí có thể tích 4 lít ở 310 $K$, đun đẳng áp đến 380 K. Thể tích sau gần bằng — phát biểu đúng là: 4,9 lít.}
{Khí có thể tích 4 lít ở 310 $K$, đun đẳng áp đến 380 K. Thể tích sau gần bằng — phát biểu 3,26 lít là đúng.}
{\True Khí đẳng áp tăng nhiệt độ tuyệt đối từ 320 $K$ lên 400 K. Tỉ số $V_2/V_1$ bằng — phát biểu đúng là: 1,25.}
{Khí đẳng áp tăng nhiệt độ tuyệt đối từ 320 $K$ lên 400 K. Tỉ số $V_2/V_1$ bằng — phát biểu 1 là đúng.}
\loigiai{
\begin{itemchoice}
\item (Đúng) Khí có thể tích 4 lít ở 310 $K$, đun đẳng áp đến 380 K. Thể tích sau gần bằng — phát biểu đúng là: 4,9 lít. (Vì): $V_2=V_1T_2/T_1=4\cdot380/310=4,9$ lít. b) Sai: trái với hệ thức hoặc điều kiện đã nêu. c) Đúng: $V_2/V_1=T_2/T_1=1,25$. d) Sai: trái với hệ thức hoặc điều kiện đã nêu
\item (Sai) Khí có thể tích 4 lít ở 310 $K$, đun đẳng áp đến 380 K. Thể tích sau gần bằng — phát biểu 3,26 lít là đúng.
\item (Đúng) Khí đẳng áp tăng nhiệt độ tuyệt đối từ 320 $K$ lên 400 K. Tỉ số $V_2/V_1$ bằng — phát biểu đúng là: 1,25.
\item (Sai) Khí đẳng áp tăng nhiệt độ tuyệt đối từ 320 $K$ lên 400 K. Tỉ số $V_2/V_1$ bằng — phát biểu 1 là đúng.
\end{itemchoice}
}
\end{ex}

% ===== Câu 14 =====
% ID: L12C2B10-04-TL
% Mức: VD
\begin{ex}
% ID: L12C2B10-04-TL
% Mức: VD
Trong dạng “Bài toán tổng hợp về quá trình đẳng áp”, 2
\choiceTF

\loigiai{
Do $V/T$ không đổi nên $T_2/T_1=V_2/V_1=2$.
}
\end{ex}

% ===== Câu 15 =====
% ID: L12C2B10-04-TLN
% Mức: VD
\begin{ex}
% ID: L12C2B10-04-TLN
% Mức: VD
Trong dạng “Bài toán tổng hợp về quá trình đẳng áp”, Thể tích đẳng áp tăng 2 lần. Nhiệt độ tuyệt đối tăng bao nhiêu lần?
\shortans{2}
\loigiai{
Do $V/T$ không đổi nên $T_2/T_1=V_2/V_1=2$.
}
\end{ex}

% ===== Câu 16 =====
% ID: L12C2B10-04-TN
% Mức: TH
\begin{ex}
% ID: L12C2B10-04-TN
% Mức: TH
Trong dạng “Bài toán tổng hợp về quá trình đẳng áp”, Khí đẳng áp tăng nhiệt độ tuyệt đối từ 320 $K$ lên 400 K. Tỉ số $V_2/V_1$ bằng
\choice
{\True 1,25}
{0,8}
{1}
{0,25}
\loigiai{
Đáp án A. $V_2/V_1=T_2/T_1=1,25$.
}
\end{ex}

% ===== Câu 17 =====
% ID: L12C2B10-05-DS
% Mức: VD
\begin{ex}
% ID: L12C2B10-05-DS
% Mức: VD
Xét các phát biểu sau trong dạng “Bài toán tổng hợp về quá trình đẳng áp”.
\choiceTF
{\True Đồ thị V- $T$ của quá trình đẳng áp trong thang kelvin là — phát biểu đúng là: đường thẳng qua gốc.}
{Đồ thị V- $T$ của quá trình đẳng áp trong thang kelvin là — phát biểu hypebol là đúng.}
{\True Khi áp suất không đổi và nhiệt độ tuyệt đối tăng gấp đôi, khối lượng riêng — phát biểu đúng là: giảm một nửa.}
{Khi áp suất không đổi và nhiệt độ tuyệt đối tăng gấp đôi, khối lượng riêng — phát biểu không đổi là đúng.}
\loigiai{
\begin{itemchoice}
\item (Đúng) Đồ thị V- $T$ của quá trình đẳng áp trong thang kelvin là — phát biểu đúng là: đường thẳng qua gốc. (Vì): $V$ tỉ lệ thuận $T$ tuyệt đối. b) Sai: trái với hệ thức hoặc điều kiện đã nêu. c) Đúng: $\rho=m/V$ và $V\sim T$, nên $\rho\sim1/T$. d) Sai: trái với hệ thức hoặc điều kiện đã nêu
\item (Sai) Đồ thị V- $T$ của quá trình đẳng áp trong thang kelvin là — phát biểu hypebol là đúng.
\item (Đúng) Khi áp suất không đổi và nhiệt độ tuyệt đối tăng gấp đôi, khối lượng riêng — phát biểu đúng là: giảm một nửa.
\item (Sai) Khi áp suất không đổi và nhiệt độ tuyệt đối tăng gấp đôi, khối lượng riêng — phát biểu không đổi là đúng.
\end{itemchoice}
}
\end{ex}

% ===== Câu 18 =====
% ID: L12C2B10-05-TL
% Mức: VD
\begin{bt}
% ID: L12C2B10-05-TL
% Mức: VD
Trong dạng “Bài toán tổng hợp về quá trình đẳng áp”, $\rho\sim1/T$, nên khối lượng riêng giảm 3 lần.
\loigiai{
$\rho\sim1/T$, nên khối lượng riêng giảm 3 lần.
}
\end{bt}

% ===== Câu 19 =====
% ID: L12C2B10-05-TLN
% Mức: VD
\begin{ex}
% ID: L12C2B10-05-TLN
% Mức: VD
Trong dạng “Bài toán tổng hợp về quá trình đẳng áp”, Nhiệt độ tuyệt đối tăng 3 lần ở áp suất không đổi. Khối lượng riêng giảm bao nhiêu lần?
\shortans{3}
\loigiai{
$\rho\sim1/T$, nên khối lượng riêng giảm 3 lần.
}
\end{ex}

% ===== Câu 20 =====
% ID: L12C2B10-05-TN
% Mức: TH
\begin{ex}
% ID: L12C2B10-05-TN
% Mức: TH
Trong dạng “Bài toán tổng hợp về quá trình đẳng áp”, Đồ thị V- $T$ của quá trình đẳng áp trong thang kelvin là
\choice
{\True đường thẳng qua gốc}
{hypebol}
{đường ngang}
{đường thẳng không qua gốc}
\loigiai{
Đáp án A. $V$ tỉ lệ thuận $T$ tuyệt đối.
}
\end{ex}

% ===== Câu 21 =====
% ID: L12C2B10-06-DS
% Mức: VD
\begin{ex}
% ID: L12C2B10-06-DS
% Mức: VD
Xét các phát biểu sau trong dạng “Bài toán tổng hợp về quá trình đẳng áp”.
\choiceTF
{\True Trong định luật Charles, nhiệt độ phải đo theo — phát biểu đúng là: kelvin.}
{Trong định luật Charles, nhiệt độ phải đo theo — phát biểu độ $C$ tùy ý là đúng.}
{\True Quả bóng mềm căng hơn khi để nơi ấm, nếu áp suất gần không đổi, chủ yếu vì — phát biểu đúng là: thể tích khí tăng theo nhiệt độ.}
{Quả bóng mềm căng hơn khi để nơi ấm, nếu áp suất gần không đổi, chủ yếu vì — phát biểu khối lượng khí tăng là đúng.}
\loigiai{
\begin{itemchoice}
\item (Đúng) Trong định luật Charles, nhiệt độ phải đo theo — phát biểu đúng là: kelvin. (Vì): Tỉ lệ thuận chỉ đúng trực tiếp với nhiệt độ tuyệt đối. b) Sai: trái với hệ thức hoặc điều kiện đã nêu. c) Đúng: Theo Charles, khi $T$ tăng thì $V$ tăng ở áp suất không đổi. d) Sai: trái với hệ thức hoặc điều kiện đã nêu
\item (Sai) Trong định luật Charles, nhiệt độ phải đo theo — phát biểu độ $C$ tùy ý là đúng.
\item (Đúng) Quả bóng mềm căng hơn khi để nơi ấm, nếu áp suất gần không đổi, chủ yếu vì — phát biểu đúng là: thể tích khí tăng theo nhiệt độ.
\item (Sai) Quả bóng mềm căng hơn khi để nơi ấm, nếu áp suất gần không đổi, chủ yếu vì — phát biểu khối lượng khí tăng là đúng.
\end{itemchoice}
}
\end{ex}

% ===== Câu 22 =====
% ID: L12C2B10-06-TL
% Mức: VDC
\begin{bt}
% ID: L12C2B10-06-TL
% Mức: VDC
Trong dạng “Bài toán tổng hợp về quá trình đẳng áp”, undefined
\loigiai{
$V_2=28\cdot300/360=23,33$ lít.
}
\end{bt}

% ===== Câu 23 =====
% ID: L12C2B10-06-TLN
% Mức: VDC
\begin{ex}
% ID: L12C2B10-06-TLN
% Mức: VDC
Trong dạng “Bài toán tổng hợp về quá trình đẳng áp”, Khí 28 lít ở 360 $K$ được làm lạnh đẳng áp về 300 K. Tính thể tích theo lít.
\shortans{23,33}
\loigiai{
$V_2=28\cdot300/360=23,33$ lít.
}
\end{ex}

% ===== Câu 24 =====
% ID: L12C2B10-06-TN
% Mức: VD
\begin{ex}
% ID: L12C2B10-06-TN
% Mức: VD
Trong dạng “Bài toán tổng hợp về quá trình đẳng áp”, Khi áp suất không đổi và nhiệt độ tuyệt đối tăng gấp đôi, khối lượng riêng
\choice
{\True giảm một nửa}
{tăng gấp đôi}
{không đổi}
{tăng bốn lần}
\loigiai{
Đáp án A. $\rho=m/V$ và $V\sim T$, nên $\rho\sim1/T$.
}
\end{ex}

% ===== Câu 25 =====
% ID: L12C2B10-07-DS
% Mức: VDC
\begin{ex}
% ID: L12C2B10-07-DS
% Mức: VDC
Xét các phát biểu sau trong dạng “Bài toán tổng hợp về quá trình đẳng áp”.
\choiceTF
{\True Đường đẳng áp có áp suất nhỏ hơn trên đồ thị V- $T$ của cùng lượng khí nằm — phát biểu đúng là: dốc hơn.}
{Đường đẳng áp có áp suất nhỏ hơn trên đồ thị V- $T$ của cùng lượng khí nằm — phát biểu thoải hơn là đúng.}
{\True Ở áp suất không đổi, thể tích ngoại suy bằng không tại — phát biểu đúng là: 0 K.}
{Ở áp suất không đổi, thể tích ngoại suy bằng không tại — phát biểu 100 ° $C$ là đúng.}
\loigiai{
\begin{itemchoice}
\item (Đúng) Đường đẳng áp có áp suất nhỏ hơn trên đồ thị V- $T$ của cùng lượng khí nằm — phát biểu đúng là: dốc hơn. (Vì): Từ $V=nRT/p$, p nhỏ hơn cho hệ số góc lớn hơn. b) Sai: trái với hệ thức hoặc điều kiện đã nêu. c) Đúng: 0 $K$ là độ không tuyệt đối theo mô hình khí lí tưởng. d) Sai: trái với hệ thức hoặc điều kiện đã nêu
\item (Sai) Đường đẳng áp có áp suất nhỏ hơn trên đồ thị V- $T$ của cùng lượng khí nằm — phát biểu thoải hơn là đúng.
\item (Đúng) Ở áp suất không đổi, thể tích ngoại suy bằng không tại — phát biểu đúng là: 0 K.
\item (Sai) Ở áp suất không đổi, thể tích ngoại suy bằng không tại — phát biểu 100 ° $C$ là đúng.
\end{itemchoice}
}
\end{ex}

% ===== Câu 26 =====
% ID: L12C2B10-07-TL
% Mức: TH
\dangbt{Khối lượng riêng của chất khí trong quá trình đẳng áp}
\begin{ex}
% ID: L12C2B10-07-TL
% Mức: TH
Trong dạng “Khối lượng riêng của chất khí trong quá trình đẳng áp”, Trình bày cơ sở lí thuyết của dạng “Khối lượng riêng của chất khí trong quá trình đẳng áp” và nêu điều kiện áp dụng.
\choiceTF

\loigiai{
Nêu đúng đại lượng không đổi, hệ thức đặc trưng và yêu cầu dùng nhiệt độ kelvin, áp suất tuyệt đối khi có liên quan.
}
\end{ex}

% ===== Câu 27 =====
% ID: L12C2B10-07-TLN
% Mức: VD
\dangbt{Cột khí bị giam bởi giọt chất lỏng trong ống tiết diện đều}
\begin{ex}
% ID: L12C2B10-07-TLN
% Mức: VD
Một ống thủy tinh tiết diện đều, một đầu kín một đầu hở, đặt thẳng đứng đầu kín ở trên. Cột khí bị giam trong ống có chiều dài $l_1 = 20$ cm ở nhiệt độ $T_1 = 300K$, được ngăn cách với khí quyển bởi một giọt thủy ngân dài $h = 5$ cm. Áp suất khí quyển $p_0 = 75$ cmHg. Hỏi khi nhiệt độ tăng lên đến $T_2 = 360K$ thì chiều dài cột khí là bao nhiêu cm?
\shortans{24 cm}
\loigiai{
Ống đặt thẳng đứng, đầu kín ở trên, đầu hở ở dưới. Giọt thủy ngân nằm phía dưới cột khí. Áp suất cột khí bị giam: $p_1 = p_0 - h = 75 - 5 = 70$ cmHg (vì đầu kín ở trên, cột khí ở trên, thủy ngân ở dưới nên áp suất khí = áp suất khí quyển trừ cột thủy ngân). Vì ống tiết diện đều và giọt thủy ngân không tràn ra ngoài, áp suất khí không đổi khi nhiệt độ thay đổi (đẳng áp): $p_2 = p_1 = 70$ cmHg. Áp dụng định luật Charles (đẳng áp): $\frac{l_1}{T_1} = \frac{l_2}{T_2}$. Suy ra: $l_2 = l_1 \cdot \frac{T_2}{T_1} = 20 \cdot \frac{360}{300} = 24$ cm.
}
\end{ex}

% ===== Câu 28 =====
% ID: L12C2B10-07-TN
% Mức: VD
\dangbt{Bài toán tổng hợp về quá trình đẳng áp}
\begin{ex}
% ID: L12C2B10-07-TN
% Mức: VD
Trong dạng “Bài toán tổng hợp về quá trình đẳng áp”, Trong định luật Charles, nhiệt độ phải đo theo
\choice
{\True kelvin}
{độ $C$ tùy ý}
{độ $F$ tùy ý}
{mọi thang đều được}
\loigiai{
Đáp án A. Tỉ lệ thuận chỉ đúng trực tiếp với nhiệt độ tuyệt đối.
}
\end{ex}

% ===== Câu 29 =====
% ID: L12C2B10-08-DS
% Mức: TH
\dangbt{Khối lượng riêng của chất khí trong quá trình đẳng áp}
\begin{ex}
% ID: L12C2B10-08-DS
% Mức: TH
Xét các phát biểu sau trong dạng “Khối lượng riêng của chất khí trong quá trình đẳng áp”.
\choiceTF
{\True Trong quá trình đẳng áp của một lượng khí xác định, đại lượng không đổi là — phát biểu đúng là: áp suất.}
{Trong quá trình đẳng áp của một lượng khí xác định, đại lượng không đổi là — phát biểu thể tích là đúng.}
{\True Hệ thức đúng của định luật Charles là — phát biểu đúng là: $V_1/T_1=V_2/T_2$.}
{Hệ thức đúng của định luật Charles là — phát biểu $p_1/T_1=p_2/T_2$ là đúng.}
\loigiai{
\begin{itemchoice}
\item (Đúng) Trong quá trình đẳng áp của một lượng khí xác định, đại lượng không đổi là — phát biểu đúng là: áp suất. (Vì): Quá trình đẳng áp có áp suất không đổi. b) Sai: trái với hệ thức hoặc điều kiện đã nêu. c) Đúng: Ở áp suất không đổi, $V/T$ là hằng số. d) Sai: trái với hệ thức hoặc điều kiện đã nêu
\item (Sai) Trong quá trình đẳng áp của một lượng khí xác định, đại lượng không đổi là — phát biểu thể tích là đúng.
\item (Đúng) Hệ thức đúng của định luật Charles là — phát biểu đúng là: $V_1/T_1=V_2/T_2$.
\item (Sai) Hệ thức đúng của định luật Charles là — phát biểu $p_1/T_1=p_2/T_2$ là đúng.
\end{itemchoice}
}
\end{ex}

% ===== Câu 30 =====
% ID: L12C2B10-08-TL
% Mức: TH
\begin{ex}
% ID: L12C2B10-08-TL
% Mức: TH
Trong dạng “Khối lượng riêng của chất khí trong quá trình đẳng áp”, Mô tả dạng đồ thị đặc trưng và cách nhận biết quá trình trong dạng “Khối lượng riêng của chất khí trong quá trình đẳng áp”.
\choiceTF

\loigiai{
Xác định các trục, đại lượng không đổi và suy ra hình dạng đồ thị từ hệ thức vật lí tương ứng.
}
\end{ex}

% ===== Câu 31 =====
% ID: L12C2B10-08-TLN
% Mức: VD
\dangbt{Cột khí bị giam bởi giọt chất lỏng trong ống tiết diện đều}
\begin{ex}
% ID: L12C2B10-08-TLN
% Mức: VD
Một ống thủy tinh tiết diện đều, hai đầu hở, đặt nằm ngang. Một giọt thủy ngân dài $h = 8$ cm nằm ở giữa ống, chia ống thành hai phần bằng nhau, mỗi phần dài $l_1 = 16$ cm ở nhiệt độ $T_1 = 280$ K. Người ta đốt nóng một đầu ống sao cho nhiệt độ cột khí ở đầu đó tăng lên $T_2 = 350K$, còn đầu kia vẫn giữ nguyên nhiệt độ $T_1 = 280$ K. Tính chiều dài cột khí ở đầu được đốt nóng theo đơn vị cm.
\shortans{20 cm}
\loigiai{
Ống thủy tinh tiết diện đều, hai đầu hở, đặt nằm ngang. Giọt thủy ngân dài $h = 8$ cm.
Ban đầu, hai cột khí có chiều dài $l_1 = 16$ cm, nhiệt độ $T_1 = 280$ K.
Áp suất hai cột khí ban đầu bằng nhau và bằng áp suất khí quyển $p_0$.
Nhiệt độ cột khí đầu nóng tăng lên $T_2 = 350K$, đầu lạnh giữ nguyên $T_1 = 280$ K.
Giọt thủy ngân cân bằng nên áp suất hai cột khí luôn bằng nhau. Quá trình biến đổi của mỗi cột khí là đẳng áp.
Tổng chiều dài hai cột khí và giọt thủy ngân không đổi: $l_2 + l_3 + h = 2l_1 + h$.
$l_2 + l_3 = 2l_1 = 2 \times 16 = 32$ cm.
Áp dụng định luật Charles cho cột khí đầu nóng: $\frac{l_2}{T_2} = \frac{l_1}{T_1}$.
$l_2 = l_1 \frac{T_2}{T_1} = 16 \times \frac{350}{280} = 16 \times \frac{5}{4} = 20$ cm.
Chiều dài cột khí ở đầu được đốt nóng là $l_2 = 20$ cm.

$20\,\text{cm}$
}
\end{ex}

% ===== Câu 32 =====
% ID: L12C2B10-08-TN
% Mức: VD
\dangbt{Bài toán tổng hợp về quá trình đẳng áp}
\begin{ex}
% ID: L12C2B10-08-TN
% Mức: VD
Trong dạng “Bài toán tổng hợp về quá trình đẳng áp”, Quả bóng mềm căng hơn khi để nơi ấm, nếu áp suất gần không đổi, chủ yếu vì
\choice
{\True thể tích khí tăng theo nhiệt độ}
{số mol khí tăng}
{khối lượng khí tăng}
{áp suất bằng không}
\loigiai{
Đáp án A. Theo Charles, khi $T$ tăng thì $V$ tăng ở áp suất không đổi.
}
\end{ex}

% ===== Câu 33 =====
% ID: L12C2B10-09-DS
% Mức: TH
\dangbt{Khối lượng riêng của chất khí trong quá trình đẳng áp}
\begin{ex}
% ID: L12C2B10-09-DS
% Mức: TH
Xét các phát biểu sau trong dạng “Khối lượng riêng của chất khí trong quá trình đẳng áp”.
\choiceTF
{\True Khí có thể tích 2 lít ở 290 $K$, đun đẳng áp đến 360 K. Thể tích sau gần bằng — phát biểu đúng là: 2,48 lít.}
{Khí có thể tích 2 lít ở 290 $K$, đun đẳng áp đến 360 K. Thể tích sau gần bằng — phát biểu 1,61 lít là đúng.}
{\True Khí đẳng áp tăng nhiệt độ tuyệt đối từ 300 $K$ lên 380 K. Tỉ số $V_2/V_1$ bằng — phát biểu đúng là: 1,27.}
{Khí đẳng áp tăng nhiệt độ tuyệt đối từ 300 $K$ lên 380 K. Tỉ số $V_2/V_1$ bằng — phát biểu 1 là đúng.}
\loigiai{
\begin{itemchoice}
\item (Đúng) Khí có thể tích 2 lít ở 290 $K$, đun đẳng áp đến 360 K. Thể tích sau gần bằng — phát biểu đúng là: 2,48 lít. (Vì): $V_2=V_1T_2/T_1=2\cdot360/290=2,48$ lít. b) Sai: trái với hệ thức hoặc điều kiện đã nêu. c) Đúng: $V_2/V_1=T_2/T_1=1,27$. d) Sai: trái với hệ thức hoặc điều kiện đã nêu
\item (Sai) Khí có thể tích 2 lít ở 290 $K$, đun đẳng áp đến 360 K. Thể tích sau gần bằng — phát biểu 1,61 lít là đúng.
\item (Đúng) Khí đẳng áp tăng nhiệt độ tuyệt đối từ 300 $K$ lên 380 K. Tỉ số $V_2/V_1$ bằng — phát biểu đúng là: 1,27.
\item (Sai) Khí đẳng áp tăng nhiệt độ tuyệt đối từ 300 $K$ lên 380 K. Tỉ số $V_2/V_1$ bằng — phát biểu 1 là đúng.
\end{itemchoice}
}
\end{ex}

% ===== Câu 34 =====
% ID: L12C2B10-09-TL
% Mức: VD
\begin{bt}
% ID: L12C2B10-09-TL
% Mức: VD
Trong dạng “Khối lượng riêng của chất khí trong quá trình đẳng áp”, Nhiệt độ tuyệt đối tăng từ 300 $K$ lên 360 $K$ ở áp suất không đổi. Tính $V_2/V_1$.
\loigiai{
$V_2/V_1=T_2/T_1=1,2$.
}
\end{bt}

% ===== Câu 35 =====
% ID: L12C2B10-09-TLN
% Mức: TH
\begin{ex}
% ID: L12C2B10-09-TLN
% Mức: TH
Trong dạng “Khối lượng riêng của chất khí trong quá trình đẳng áp”, Khí 7 lít ở 280 $K$ được đun đẳng áp đến 340 K. Tính thể tích sau theo lít.
\shortans{8,5}
\loigiai{
$V_2=7\cdot340/280=8,5$ lít.
}
\end{ex}

% ===== Câu 36 =====
% ID: L12C2B10-09-TN
% Mức: VDC
\dangbt{Bài toán tổng hợp về quá trình đẳng áp}
\begin{ex}
% ID: L12C2B10-09-TN
% Mức: VDC
Trong dạng “Bài toán tổng hợp về quá trình đẳng áp”, Đường đẳng áp có áp suất nhỏ hơn trên đồ thị V- $T$ của cùng lượng khí nằm
\choice
{\True dốc hơn}
{thoải hơn}
{trùng trục $T$}
{luôn trùng nhau}
\loigiai{
Đáp án A. Từ $V=nRT/p$, p nhỏ hơn cho hệ số góc lớn hơn.
}
\end{ex}

% ===== Câu 37 =====
% ID: L12C2B10-10-DS
% Mức: VD
\dangbt{Khối lượng riêng của chất khí trong quá trình đẳng áp}
\begin{ex}
% ID: L12C2B10-10-DS
% Mức: VD
Xét các phát biểu sau trong dạng “Khối lượng riêng của chất khí trong quá trình đẳng áp”.
\choiceTF
{\True Đồ thị V- $T$ của quá trình đẳng áp trong thang kelvin là — phát biểu đúng là: đường thẳng qua gốc.}
{Đồ thị V- $T$ của quá trình đẳng áp trong thang kelvin là — phát biểu hypebol là đúng.}
{\True Khi áp suất không đổi và nhiệt độ tuyệt đối tăng gấp đôi, khối lượng riêng — phát biểu đúng là: giảm một nửa.}
{Khi áp suất không đổi và nhiệt độ tuyệt đối tăng gấp đôi, khối lượng riêng — phát biểu không đổi là đúng.}
\loigiai{
\begin{itemchoice}
\item (Đúng) Đồ thị V- $T$ của quá trình đẳng áp trong thang kelvin là — phát biểu đúng là: đường thẳng qua gốc. (Vì): $V$ tỉ lệ thuận $T$ tuyệt đối. b) Sai: trái với hệ thức hoặc điều kiện đã nêu. c) Đúng: $\rho=m/V$ và $V\sim T$, nên $\rho\sim1/T$. d) Sai: trái với hệ thức hoặc điều kiện đã nêu
\item (Sai) Đồ thị V- $T$ của quá trình đẳng áp trong thang kelvin là — phát biểu hypebol là đúng.
\item (Đúng) Khi áp suất không đổi và nhiệt độ tuyệt đối tăng gấp đôi, khối lượng riêng — phát biểu đúng là: giảm một nửa.
\item (Sai) Khi áp suất không đổi và nhiệt độ tuyệt đối tăng gấp đôi, khối lượng riêng — phát biểu không đổi là đúng.
\end{itemchoice}
}
\end{ex}

% ===== Câu 38 =====
% ID: L12C2B10-10-TL
% Mức: VD
\begin{ex}
% ID: L12C2B10-10-TL
% Mức: VD
Trong dạng “Khối lượng riêng của chất khí trong quá trình đẳng áp”, 3
\choiceTF

\loigiai{
Do $V/T$ không đổi nên $T_2/T_1=V_2/V_1=3$.
}
\end{ex}

% ===== Câu 39 =====
% ID: L12C2B10-10-TLN
% Mức: TH
\begin{ex}
% ID: L12C2B10-10-TLN
% Mức: TH
Trong dạng “Khối lượng riêng của chất khí trong quá trình đẳng áp”, Khí có thể tích 3 lít ở 290 $K$, đẳng áp tăng thể tích lên 12 lít. Tính nhiệt độ sau theo K.
\shortans{1160}
\loigiai{
$T_2=T_1V_2/V_1=290\cdot4=1160$ K.
}
\end{ex}

% ===== Câu 40 =====
% ID: L12C2B10-10-TN
% Mức: VDC
\dangbt{Bài toán tổng hợp về quá trình đẳng áp}
\begin{ex}
% ID: L12C2B10-10-TN
% Mức: VDC
Trong dạng “Bài toán tổng hợp về quá trình đẳng áp”, Ở áp suất không đổi, thể tích ngoại suy bằng không tại
\choice
{\True 0 $K$}
{0 ° $C$}
{100 ° $C$}
{273 $K$}
\loigiai{
Đáp án A. 0 $K$ là độ không tuyệt đối theo mô hình khí lí tưởng.
}
\end{ex}

% ===== Câu 41 =====
% ID: L12C2B10-11-DS
% Mức: VD
\dangbt{Khối lượng riêng của chất khí trong quá trình đẳng áp}
\begin{ex}
% ID: L12C2B10-11-DS
% Mức: VD
Xét các phát biểu sau trong dạng “Khối lượng riêng của chất khí trong quá trình đẳng áp”.
\choiceTF
{\True Trong định luật Charles, nhiệt độ phải đo theo — phát biểu đúng là: kelvin.}
{Trong định luật Charles, nhiệt độ phải đo theo — phát biểu độ $C$ tùy ý là đúng.}
{\True Quả bóng mềm căng hơn khi để nơi ấm, nếu áp suất gần không đổi, chủ yếu vì — phát biểu đúng là: thể tích khí tăng theo nhiệt độ.}
{Quả bóng mềm căng hơn khi để nơi ấm, nếu áp suất gần không đổi, chủ yếu vì — phát biểu khối lượng khí tăng là đúng.}
\loigiai{
\begin{itemchoice}
\item (Đúng) Trong định luật Charles, nhiệt độ phải đo theo — phát biểu đúng là: kelvin. (Vì): Tỉ lệ thuận chỉ đúng trực tiếp với nhiệt độ tuyệt đối. b) Sai: trái với hệ thức hoặc điều kiện đã nêu. c) Đúng: Theo Charles, khi $T$ tăng thì $V$ tăng ở áp suất không đổi. d) Sai: trái với hệ thức hoặc điều kiện đã nêu
\item (Sai) Trong định luật Charles, nhiệt độ phải đo theo — phát biểu độ $C$ tùy ý là đúng.
\item (Đúng) Quả bóng mềm căng hơn khi để nơi ấm, nếu áp suất gần không đổi, chủ yếu vì — phát biểu đúng là: thể tích khí tăng theo nhiệt độ.
\item (Sai) Quả bóng mềm căng hơn khi để nơi ấm, nếu áp suất gần không đổi, chủ yếu vì — phát biểu khối lượng khí tăng là đúng.
\end{itemchoice}
}
\end{ex}

% ===== Câu 42 =====
% ID: L12C2B10-11-TL
% Mức: VD
\begin{bt}
% ID: L12C2B10-11-TL
% Mức: VD
Trong dạng “Khối lượng riêng của chất khí trong quá trình đẳng áp”, $\rho\sim1/T$, nên khối lượng riêng giảm 4 lần.
\loigiai{
$\rho\sim1/T$, nên khối lượng riêng giảm 4 lần.
}
\end{bt}

% ===== Câu 43 =====
% ID: L12C2B10-11-TLN
% Mức: VD
\begin{ex}
% ID: L12C2B10-11-TLN
% Mức: VD
Trong dạng “Khối lượng riêng của chất khí trong quá trình đẳng áp”, Nhiệt độ tuyệt đối tăng từ 300 $K$ lên 360 $K$ ở áp suất không đổi. Tính $V_2/V_1$.
\shortans{1,2}
\loigiai{
$V_2/V_1=T_2/T_1=1,2$.
}
\end{ex}

% ===== Câu 44 =====
% ID: L12C2B10-11-TN
% Mức: VD
\dangbt{Cột khí bị giam bởi giọt chất lỏng trong ống tiết diện đều}
\begin{ex}
% ID: L12C2B10-11-TN
% Mức: VD
Một ống thủy tinh tiết diện đều nằm ngang, một đầu kín, giam cột khí dài $L_1 = 24\ \mathrm{cm}$ ở $27\ \mathrm{^\circ C}$ nhờ một giọt dầu. Làm lạnh ống xuống $-23\ \mathrm{^\circ C}$. Cột khí ngắn đi bao nhiêu xentimét?
\choice
{$3\ \mathrm{cm}$}
{\True $4\ \mathrm{cm}$}
{$6\ \mathrm{cm}$}
{$8\ \mathrm{cm}$}
\loigiai{
Đổi nhiệt độ: $T_1 = 300\ \mathrm{K}$, $T_2 = 250\ \mathrm{K}$. Ống nằm ngang, áp suất không đổi, áp dụng định luật Charles: $L_2 = L_1 \cdot \frac{T_2}{T_1} = 24 \cdot \frac{250}{300} = 20\ \mathrm{cm}$. Cột khí ngắn đi: $\Delta L = 24 - 20 = 4\ \mathrm{cm}$. Đáp án đúng là B.
}
\end{ex}

% ===== Câu 45 =====
% ID: L12C2B10-12-DS
% Mức: VDC
\dangbt{Khối lượng riêng của chất khí trong quá trình đẳng áp}
\begin{ex}
% ID: L12C2B10-12-DS
% Mức: VDC
Xét các phát biểu sau trong dạng “Khối lượng riêng của chất khí trong quá trình đẳng áp”.
\choiceTF
{\True Đường đẳng áp có áp suất nhỏ hơn trên đồ thị V- $T$ của cùng lượng khí nằm — phát biểu đúng là: dốc hơn.}
{Đường đẳng áp có áp suất nhỏ hơn trên đồ thị V- $T$ của cùng lượng khí nằm — phát biểu thoải hơn là đúng.}
{\True Ở áp suất không đổi, thể tích ngoại suy bằng không tại — phát biểu đúng là: 0 K.}
{Ở áp suất không đổi, thể tích ngoại suy bằng không tại — phát biểu 100 ° $C$ là đúng.}
\loigiai{
\begin{itemchoice}
\item (Đúng) Đường đẳng áp có áp suất nhỏ hơn trên đồ thị V- $T$ của cùng lượng khí nằm — phát biểu đúng là: dốc hơn. (Vì): Từ $V=nRT/p$, p nhỏ hơn cho hệ số góc lớn hơn. b) Sai: trái với hệ thức hoặc điều kiện đã nêu. c) Đúng: 0 $K$ là độ không tuyệt đối theo mô hình khí lí tưởng. d) Sai: trái với hệ thức hoặc điều kiện đã nêu
\item (Sai) Đường đẳng áp có áp suất nhỏ hơn trên đồ thị V- $T$ của cùng lượng khí nằm — phát biểu thoải hơn là đúng.
\item (Đúng) Ở áp suất không đổi, thể tích ngoại suy bằng không tại — phát biểu đúng là: 0 K.
\item (Sai) Ở áp suất không đổi, thể tích ngoại suy bằng không tại — phát biểu 100 ° $C$ là đúng.
\end{itemchoice}
}
\end{ex}

% ===== Câu 46 =====
% ID: L12C2B10-12-TL
% Mức: VDC
\begin{bt}
% ID: L12C2B10-12-TL
% Mức: VDC
Trong dạng “Khối lượng riêng của chất khí trong quá trình đẳng áp”, undefined
\loigiai{
$V_2=10\cdot280/340=8,24$ lít.
}
\end{bt}

% ===== Câu 47 =====
% ID: L12C2B10-12-TLN
% Mức: VD
\begin{ex}
% ID: L12C2B10-12-TLN
% Mức: VD
Trong dạng “Khối lượng riêng của chất khí trong quá trình đẳng áp”, Thể tích đẳng áp tăng 3 lần. Nhiệt độ tuyệt đối tăng bao nhiêu lần?
\shortans{3}
\loigiai{
Do $V/T$ không đổi nên $T_2/T_1=V_2/V_1=3$.
}
\end{ex}

% ===== Câu 48 =====
% ID: L12C2B10-12-TN
% Mức: VD
\dangbt{Cột khí bị giam bởi giọt chất lỏng trong ống tiết diện đều}
\begin{ex}
% ID: L12C2B10-12-TN
% Mức: VD
Một ống thủy tinh tiết diện đều nằm ngang, một đầu kín, giam cột khí nhờ một giọt thủy ngân. Ở $17\ \mathrm{^\circ C}$ cột khí dài $20\ \mathrm{cm}$. Phải nung ống đến nhiệt độ bao nhiêu để cột khí dài $25\ \mathrm{cm}$?
\choice
{$21,25\ \mathrm{^\circ C}$}
{$363,75\ \mathrm{^\circ C}$}
{\True $90\ \mathrm{^\circ C}$}
{$72,5\ \mathrm{^\circ C}$}
\loigiai{
Đổi $T_1 = 290\ \mathrm{K}$. Ống nằm ngang, áp suất không đổi, áp dụng định luật Charles: $\frac{L_1}{T_1} = \frac{L_2}{T_2}$. Suy ra $T_2 = T_1 \cdot \frac{L_2}{L_1} = 290 \cdot \frac{25}{20} = 290 \cdot 1,25 = 362,5\ \mathrm{K}$. Đổi ra độ C: $t_2 = 362,5 - 273 = 89,5 \approx 90\ \mathrm{^\circ C}$. Đáp án đúng là C.
}
\end{ex}

% ===== Câu 49 =====
% ID: L12C2B10-13-DS
% Mức: TH
\dangbt{Nhận biết quá trình đẳng áp và định luật Charles}
\begin{ex}
% ID: L12C2B10-13-DS
% Mức: TH
Xét các phát biểu sau trong dạng “Nhận biết quá trình đẳng áp và định luật Charles”.
\choiceTF
{\True Trong quá trình đẳng áp của một lượng khí xác định, đại lượng không đổi là — phát biểu đúng là: áp suất.}
{Trong quá trình đẳng áp của một lượng khí xác định, đại lượng không đổi là — phát biểu thể tích là đúng.}
{\True Hệ thức đúng của định luật Charles là — phát biểu đúng là: $V_1/T_1=V_2/T_2$.}
{Hệ thức đúng của định luật Charles là — phát biểu $p_1/T_1=p_2/T_2$ là đúng.}
\loigiai{
\begin{itemchoice}
\item (Đúng) Trong quá trình đẳng áp của một lượng khí xác định, đại lượng không đổi là — phát biểu đúng là: áp suất. (Vì): Quá trình đẳng áp có áp suất không đổi. b) Sai: trái với hệ thức hoặc điều kiện đã nêu. c) Đúng: Ở áp suất không đổi, $V/T$ là hằng số. d) Sai: trái với hệ thức hoặc điều kiện đã nêu
\item (Sai) Trong quá trình đẳng áp của một lượng khí xác định, đại lượng không đổi là — phát biểu thể tích là đúng.
\item (Đúng) Hệ thức đúng của định luật Charles là — phát biểu đúng là: $V_1/T_1=V_2/T_2$.
\item (Sai) Hệ thức đúng của định luật Charles là — phát biểu $p_1/T_1=p_2/T_2$ là đúng.
\end{itemchoice}
}
\end{ex}

% ===== Câu 50 =====
% ID: L12C2B10-13-TL
% Mức: TH
\begin{ex}
% ID: L12C2B10-13-TL
% Mức: TH
Trong dạng “Nhận biết quá trình đẳng áp và định luật Charles”, Trình bày cơ sở lí thuyết của dạng “Nhận biết quá trình đẳng áp và định luật Charles” và nêu điều kiện áp dụng.
\choiceTF

\loigiai{
Nêu đúng đại lượng không đổi, hệ thức đặc trưng và yêu cầu dùng nhiệt độ kelvin, áp suất tuyệt đối khi có liên quan.
}
\end{ex}

% ===== Câu 51 =====
% ID: L12C2B10-13-TLN
% Mức: VD
\dangbt{Khối lượng riêng của chất khí trong quá trình đẳng áp}
\begin{ex}
% ID: L12C2B10-13-TLN
% Mức: VD
Trong dạng “Khối lượng riêng của chất khí trong quá trình đẳng áp”, Nhiệt độ tuyệt đối tăng 4 lần ở áp suất không đổi. Khối lượng riêng giảm bao nhiêu lần?
\shortans{4}
\loigiai{
$\rho\sim1/T$, nên khối lượng riêng giảm 4 lần.
}
\end{ex}

% ===== Câu 52 =====
% ID: L12C2B10-13-TN
% Mức: VD
\dangbt{Cột khí bị giam bởi giọt chất lỏng trong ống tiết diện đều}
\begin{ex}
% ID: L12C2B10-13-TN
% Mức: VD
Một ống thủy tinh tiết diện đều nằm ngang, một đầu kín một đầu hở, giam một cột khí dài $L_1 = 15\ \mathrm{cm}$ ở nhiệt độ $T_1 = 300\ \mathrm{K}$ nhờ một giọt thủy ngân. Nung nóng ống đến nhiệt độ $T_2 = 360\ \mathrm{K}$. Chiều dài cột khí lúc sau là bao nhiêu?
\choice
{\True $18\ \mathrm{cm}$}
{$15\ \mathrm{cm}$}
{$20\ \mathrm{cm}$}
{$12\ \mathrm{cm}$}
\loigiai{
Ống nằm ngang, giọt thủy ngân dịch chuyển tự do nên áp suất khí không đổi. Áp dụng định luật Charles: $\frac{L_1}{T_1} = \frac{L_2}{T_2}$. Suy ra $L_2 = L_1 \cdot \frac{T_2}{T_1} = 15 \cdot \frac{360}{300} = 15 \cdot 1,2 = 18\ \mathrm{cm}$. Đáp án đúng là A.
}
\end{ex}

% ===== Câu 53 =====
% ID: L12C2B10-14-DS
% Mức: TH
\dangbt{Nhận biết quá trình đẳng áp và định luật Charles}
\begin{ex}
% ID: L12C2B10-14-DS
% Mức: TH
Xét các phát biểu sau trong dạng “Nhận biết quá trình đẳng áp và định luật Charles”.
\choiceTF
{\True Khí có thể tích 6 lít ở 260 $K$, đun đẳng áp đến 330 K. Thể tích sau gần bằng — phát biểu đúng là: 7,62 lít.}
{Khí có thể tích 6 lít ở 260 $K$, đun đẳng áp đến 330 K. Thể tích sau gần bằng — phát biểu 4,73 lít là đúng.}
{\True Khí đẳng áp tăng nhiệt độ tuyệt đối từ 270 $K$ lên 350 K. Tỉ số $V_2/V_1$ bằng — phát biểu đúng là: 1,3.}
{Khí đẳng áp tăng nhiệt độ tuyệt đối từ 270 $K$ lên 350 K. Tỉ số $V_2/V_1$ bằng — phát biểu 1 là đúng.}
\loigiai{
\begin{itemchoice}
\item (Đúng) Khí có thể tích 6 lít ở 260 $K$, đun đẳng áp đến 330 K. Thể tích sau gần bằng — phát biểu đúng là: 7,62 lít. (Vì): $V_2=V_1T_2/T_1=6\cdot330/260=7,62$ lít. b) Sai: trái với hệ thức hoặc điều kiện đã nêu. c) Đúng: $V_2/V_1=T_2/T_1=1,3$. d) Sai: trái với hệ thức hoặc điều kiện đã nêu
\item (Sai) Khí có thể tích 6 lít ở 260 $K$, đun đẳng áp đến 330 K. Thể tích sau gần bằng — phát biểu 4,73 lít là đúng.
\item (Đúng) Khí đẳng áp tăng nhiệt độ tuyệt đối từ 270 $K$ lên 350 K. Tỉ số $V_2/V_1$ bằng — phát biểu đúng là: 1,3.
\item (Sai) Khí đẳng áp tăng nhiệt độ tuyệt đối từ 270 $K$ lên 350 K. Tỉ số $V_2/V_1$ bằng — phát biểu 1 là đúng.
\end{itemchoice}
}
\end{ex}

% ===== Câu 54 =====
% ID: L12C2B10-14-TL
% Mức: TH
\begin{ex}
% ID: L12C2B10-14-TL
% Mức: TH
Trong dạng “Nhận biết quá trình đẳng áp và định luật Charles”, Mô tả dạng đồ thị đặc trưng và cách nhận biết quá trình trong dạng “Nhận biết quá trình đẳng áp và định luật Charles”.
\choiceTF

\loigiai{
Xác định các trục, đại lượng không đổi và suy ra hình dạng đồ thị từ hệ thức vật lí tương ứng.
}
\end{ex}

% ===== Câu 55 =====
% ID: L12C2B10-14-TLN
% Mức: VDC
\dangbt{Khối lượng riêng của chất khí trong quá trình đẳng áp}
\begin{ex}
% ID: L12C2B10-14-TLN
% Mức: VDC
Trong dạng “Khối lượng riêng của chất khí trong quá trình đẳng áp”, Khí 10 lít ở 340 $K$ được làm lạnh đẳng áp về 280 K. Tính thể tích theo lít.
\shortans{8,24}
\loigiai{
$V_2=10\cdot280/340=8,24$ lít.
}
\end{ex}

% ===== Câu 56 =====
% ID: L12C2B10-14-TN
% Mức: VD
\dangbt{Cột khí bị giam bởi giọt chất lỏng trong ống tiết diện đều}
\begin{ex}
% ID: L12C2B10-14-TN
% Mức: VD
Một ống thủy tinh tiết diện đều nằm ngang, một đầu kín, giam cột khí dài $L_1$ ở $7\ \mathrm{^\circ C}$ nhờ một giọt thủy ngân. Khi nung ống đến $77\ \mathrm{^\circ C}$, cột khí tăng thêm $7\ \mathrm{cm}$. Chiều dài ban đầu $L_1$ của cột khí là bao nhiêu?
\choice
{$21\ \mathrm{cm}$}
{$28\ \mathrm{cm}$}
{$35\ \mathrm{cm}$}
{\True $28\ \mathrm{cm}$}
\loigiai{
Đổi nhiệt độ: $T_1 = 280\ \mathrm{K}$, $T_2 = 350\ \mathrm{K}$. Ống nằm ngang, áp suất không đổi, áp dụng định luật Charles: $\frac{L_1}{T_1} = \frac{L_2}{T_2}$, với $L_2 = L_1 + 7$. Ta có: $\frac{L_1}{280} = \frac{L_1 + 7}{350}$. Suy ra $350 L_1 = 280(L_1 + 7) = 280 L_1 + 1960$. Do đó $70 L_1 = 1960$, $L_1 = 28\ \mathrm{cm}$. Đáp án đúng là D.
}
\end{ex}

% ===== Câu 57 =====
% ID: L12C2B10-15-DS
% Mức: VD
\dangbt{Nhận biết quá trình đẳng áp và định luật Charles}
\begin{ex}
% ID: L12C2B10-15-DS
% Mức: VD
Xét các phát biểu sau trong dạng “Nhận biết quá trình đẳng áp và định luật Charles”.
\choiceTF
{\True Đồ thị V- $T$ của quá trình đẳng áp trong thang kelvin là — phát biểu đúng là: đường thẳng qua gốc.}
{Đồ thị V- $T$ của quá trình đẳng áp trong thang kelvin là — phát biểu hypebol là đúng.}
{\True Khi áp suất không đổi và nhiệt độ tuyệt đối tăng gấp đôi, khối lượng riêng — phát biểu đúng là: giảm một nửa.}
{Khi áp suất không đổi và nhiệt độ tuyệt đối tăng gấp đôi, khối lượng riêng — phát biểu không đổi là đúng.}
\loigiai{
\begin{itemchoice}
\item (Đúng) Đồ thị V- $T$ của quá trình đẳng áp trong thang kelvin là — phát biểu đúng là: đường thẳng qua gốc. (Vì): $V$ tỉ lệ thuận $T$ tuyệt đối. b) Sai: trái với hệ thức hoặc điều kiện đã nêu. c) Đúng: $\rho=m/V$ và $V\sim T$, nên $\rho\sim1/T$. d) Sai: trái với hệ thức hoặc điều kiện đã nêu
\item (Sai) Đồ thị V- $T$ của quá trình đẳng áp trong thang kelvin là — phát biểu hypebol là đúng.
\item (Đúng) Khi áp suất không đổi và nhiệt độ tuyệt đối tăng gấp đôi, khối lượng riêng — phát biểu đúng là: giảm một nửa.
\item (Sai) Khi áp suất không đổi và nhiệt độ tuyệt đối tăng gấp đôi, khối lượng riêng — phát biểu không đổi là đúng.
\end{itemchoice}
}
\end{ex}

% ===== Câu 58 =====
% ID: L12C2B10-15-TL
% Mức: VD
\begin{bt}
% ID: L12C2B10-15-TL
% Mức: VD
Trong dạng “Nhận biết quá trình đẳng áp và định luật Charles”, Nhiệt độ tuyệt đối tăng từ 320 $K$ lên 380 $K$ ở áp suất không đổi. Tính $V_2/V_1$.
\loigiai{
$V_2/V_1=T_2/T_1=1,19$.
}
\end{bt}

% ===== Câu 59 =====
% ID: L12C2B10-15-TLN
% Mức: TH
\begin{ex}
% ID: L12C2B10-15-TLN
% Mức: TH
Trong dạng “Nhận biết quá trình đẳng áp và định luật Charles”, Khí 4 lít ở 300 $K$ được đun đẳng áp đến 360 K. Tính thể tích sau theo lít.
\shortans{4,8}
\loigiai{
$V_2=4\cdot360/300=4,8$ lít.
}
\end{ex}

% ===== Câu 60 =====
% ID: L12C2B10-15-TN
% Mức: NB
\dangbt{Khối lượng riêng của chất khí trong quá trình đẳng áp}
\begin{ex}
% ID: L12C2B10-15-TN
% Mức: NB
Trong dạng “Khối lượng riêng của chất khí trong quá trình đẳng áp”, Trong quá trình đẳng áp của một lượng khí xác định, đại lượng không đổi là
\choice
{\True áp suất}
{thể tích}
{nhiệt độ}
{khối lượng riêng}
\loigiai{
Đáp án A. Quá trình đẳng áp có áp suất không đổi.
}
\end{ex}

% ===== Câu 61 =====
% ID: L12C2B10-16-DS
% Mức: VD
\dangbt{Nhận biết quá trình đẳng áp và định luật Charles}
\begin{ex}
% ID: L12C2B10-16-DS
% Mức: VD
Xét các phát biểu sau trong dạng “Nhận biết quá trình đẳng áp và định luật Charles”.
\choiceTF
{\True Trong định luật Charles, nhiệt độ phải đo theo — phát biểu đúng là: kelvin.}
{Trong định luật Charles, nhiệt độ phải đo theo — phát biểu độ $C$ tùy ý là đúng.}
{\True Quả bóng mềm căng hơn khi để nơi ấm, nếu áp suất gần không đổi, chủ yếu vì — phát biểu đúng là: thể tích khí tăng theo nhiệt độ.}
{Quả bóng mềm căng hơn khi để nơi ấm, nếu áp suất gần không đổi, chủ yếu vì — phát biểu khối lượng khí tăng là đúng.}
\loigiai{
\begin{itemchoice}
\item (Đúng) Trong định luật Charles, nhiệt độ phải đo theo — phát biểu đúng là: kelvin. (Vì): Tỉ lệ thuận chỉ đúng trực tiếp với nhiệt độ tuyệt đối. b) Sai: trái với hệ thức hoặc điều kiện đã nêu. c) Đúng: Theo Charles, khi $T$ tăng thì $V$ tăng ở áp suất không đổi. d) Sai: trái với hệ thức hoặc điều kiện đã nêu
\item (Sai) Trong định luật Charles, nhiệt độ phải đo theo — phát biểu độ $C$ tùy ý là đúng.
\item (Đúng) Quả bóng mềm căng hơn khi để nơi ấm, nếu áp suất gần không đổi, chủ yếu vì — phát biểu đúng là: thể tích khí tăng theo nhiệt độ.
\item (Sai) Quả bóng mềm căng hơn khi để nơi ấm, nếu áp suất gần không đổi, chủ yếu vì — phát biểu khối lượng khí tăng là đúng.
\end{itemchoice}
}
\end{ex}

% ===== Câu 62 =====
% ID: L12C2B10-16-TL
% Mức: VD
\begin{ex}
% ID: L12C2B10-16-TL
% Mức: VD
Trong dạng “Nhận biết quá trình đẳng áp và định luật Charles”, 3
\choiceTF

\loigiai{
Do $V/T$ không đổi nên $T_2/T_1=V_2/V_1=3$.
}
\end{ex}

% ===== Câu 63 =====
% ID: L12C2B10-16-TLN
% Mức: TH
\begin{ex}
% ID: L12C2B10-16-TLN
% Mức: TH
Trong dạng “Nhận biết quá trình đẳng áp và định luật Charles”, Khí có thể tích 6 lít ở 310 $K$, đẳng áp tăng thể tích lên 24 lít. Tính nhiệt độ sau theo K.
\shortans{1240}
\loigiai{
$T_2=T_1V_2/V_1=310\cdot4=1240$ K.
}
\end{ex}

% ===== Câu 64 =====
% ID: L12C2B10-16-TN
% Mức: NB
\dangbt{Khối lượng riêng của chất khí trong quá trình đẳng áp}
\begin{ex}
% ID: L12C2B10-16-TN
% Mức: NB
Trong dạng “Khối lượng riêng của chất khí trong quá trình đẳng áp”, Hệ thức đúng của định luật Charles là
\choice
{\True $V_1/T_1=V_2/T_2$}
{$p_1V_1=p_2V_2$}
{$p_1/T_1=p_2/T_2$}
{$V_1T_1=V_2T_2$}
\loigiai{
Đáp án A. Ở áp suất không đổi, $V/T$ là hằng số.
}
\end{ex}

% ===== Câu 65 =====
% ID: L12C2B10-17-DS
% Mức: VDC
\dangbt{Nhận biết quá trình đẳng áp và định luật Charles}
\begin{ex}
% ID: L12C2B10-17-DS
% Mức: VDC
Xét các phát biểu sau trong dạng “Nhận biết quá trình đẳng áp và định luật Charles”.
\choiceTF
{\True Đường đẳng áp có áp suất nhỏ hơn trên đồ thị V- $T$ của cùng lượng khí nằm — phát biểu đúng là: dốc hơn.}
{Đường đẳng áp có áp suất nhỏ hơn trên đồ thị V- $T$ của cùng lượng khí nằm — phát biểu thoải hơn là đúng.}
{\True Ở áp suất không đổi, thể tích ngoại suy bằng không tại — phát biểu đúng là: 0 K.}
{Ở áp suất không đổi, thể tích ngoại suy bằng không tại — phát biểu 100 ° $C$ là đúng.}
\loigiai{
\begin{itemchoice}
\item (Đúng) Đường đẳng áp có áp suất nhỏ hơn trên đồ thị V- $T$ của cùng lượng khí nằm — phát biểu đúng là: dốc hơn. (Vì): Từ $V=nRT/p$, p nhỏ hơn cho hệ số góc lớn hơn. b) Sai: trái với hệ thức hoặc điều kiện đã nêu. c) Đúng: 0 $K$ là độ không tuyệt đối theo mô hình khí lí tưởng. d) Sai: trái với hệ thức hoặc điều kiện đã nêu
\item (Sai) Đường đẳng áp có áp suất nhỏ hơn trên đồ thị V- $T$ của cùng lượng khí nằm — phát biểu thoải hơn là đúng.
\item (Đúng) Ở áp suất không đổi, thể tích ngoại suy bằng không tại — phát biểu đúng là: 0 K.
\item (Sai) Ở áp suất không đổi, thể tích ngoại suy bằng không tại — phát biểu 100 ° $C$ là đúng.
\end{itemchoice}
}
\end{ex}

% ===== Câu 66 =====
% ID: L12C2B10-17-TL
% Mức: VD
\begin{bt}
% ID: L12C2B10-17-TL
% Mức: VD
Trong dạng “Nhận biết quá trình đẳng áp và định luật Charles”, $\rho\sim1/T$, nên khối lượng riêng giảm 4 lần.
\loigiai{
$\rho\sim1/T$, nên khối lượng riêng giảm 4 lần.
}
\end{bt}

% ===== Câu 67 =====
% ID: L12C2B10-17-TLN
% Mức: VD
\begin{ex}
% ID: L12C2B10-17-TLN
% Mức: VD
Trong dạng “Nhận biết quá trình đẳng áp và định luật Charles”, Nhiệt độ tuyệt đối tăng từ 320 $K$ lên 380 $K$ ở áp suất không đổi. Tính $V_2/V_1$.
\shortans{1,19}
\loigiai{
$V_2/V_1=T_2/T_1=1,19$.
}
\end{ex}

% ===== Câu 68 =====
% ID: L12C2B10-17-TN
% Mức: TH
\dangbt{Khối lượng riêng của chất khí trong quá trình đẳng áp}
\begin{ex}
% ID: L12C2B10-17-TN
% Mức: TH
Trong dạng “Khối lượng riêng của chất khí trong quá trình đẳng áp”, Khí có thể tích 2 lít ở 290 $K$, đun đẳng áp đến 360 K. Thể tích sau gần bằng
\choice
{\True 2,48 lít}
{1,61 lít}
{2 lít}
{4,97 lít}
\loigiai{
Đáp án A. $V_2=V_1T_2/T_1=2\cdot360/290=2,48$ lít.
}
\end{ex}

% ===== Câu 69 =====
% ID: L12C2B10-18-DS
% Mức: TH
\dangbt{Tính thể tích và nhiệt độ trong quá trình đẳng áp}
\begin{ex}
% ID: L12C2B10-18-DS
% Mức: TH
Xét các phát biểu sau trong dạng “Tính thể tích và nhiệt độ trong quá trình đẳng áp”.
\choiceTF
{\True Trong quá trình đẳng áp của một lượng khí xác định, đại lượng không đổi là — phát biểu đúng là: áp suất.}
{Trong quá trình đẳng áp của một lượng khí xác định, đại lượng không đổi là — phát biểu thể tích là đúng.}
{\True Hệ thức đúng của định luật Charles là — phát biểu đúng là: $V_1/T_1=V_2/T_2$.}
{Hệ thức đúng của định luật Charles là — phát biểu $p_1/T_1=p_2/T_2$ là đúng.}
\loigiai{
\begin{itemchoice}
\item (Đúng) Trong quá trình đẳng áp của một lượng khí xác định, đại lượng không đổi là — phát biểu đúng là: áp suất. (Vì): Quá trình đẳng áp có áp suất không đổi. b) Sai: trái với hệ thức hoặc điều kiện đã nêu. c) Đúng: Ở áp suất không đổi, $V/T$ là hằng số. d) Sai: trái với hệ thức hoặc điều kiện đã nêu
\item (Sai) Trong quá trình đẳng áp của một lượng khí xác định, đại lượng không đổi là — phát biểu thể tích là đúng.
\item (Đúng) Hệ thức đúng của định luật Charles là — phát biểu đúng là: $V_1/T_1=V_2/T_2$.
\item (Sai) Hệ thức đúng của định luật Charles là — phát biểu $p_1/T_1=p_2/T_2$ là đúng.
\end{itemchoice}
}
\end{ex}

% ===== Câu 70 =====
% ID: L12C2B10-18-TL
% Mức: VDC
\dangbt{Nhận biết quá trình đẳng áp và định luật Charles}
\begin{bt}
% ID: L12C2B10-18-TL
% Mức: VDC
Trong dạng “Nhận biết quá trình đẳng áp và định luật Charles”, undefined
\loigiai{
$V_2=16\cdot300/360=13,33$ lít.
}
\end{bt}

% ===== Câu 71 =====
% ID: L12C2B10-18-TLN
% Mức: VD
\begin{ex}
% ID: L12C2B10-18-TLN
% Mức: VD
Trong dạng “Nhận biết quá trình đẳng áp và định luật Charles”, Thể tích đẳng áp tăng 3 lần. Nhiệt độ tuyệt đối tăng bao nhiêu lần?
\shortans{3}
\loigiai{
Do $V/T$ không đổi nên $T_2/T_1=V_2/V_1=3$.
}
\end{ex}

% ===== Câu 72 =====
% ID: L12C2B10-18-TN
% Mức: TH
\dangbt{Khối lượng riêng của chất khí trong quá trình đẳng áp}
\begin{ex}
% ID: L12C2B10-18-TN
% Mức: TH
Trong dạng “Khối lượng riêng của chất khí trong quá trình đẳng áp”, Khí đẳng áp tăng nhiệt độ tuyệt đối từ 300 $K$ lên 380 K. Tỉ số $V_2/V_1$ bằng
\choice
{\True 1,27}
{0,79}
{1}
{0,27}
\loigiai{
Đáp án A. $V_2/V_1=T_2/T_1=1,27$.
}
\end{ex}

% ===== Câu 73 =====
% ID: L12C2B10-19-DS
% Mức: TH
\dangbt{Tính thể tích và nhiệt độ trong quá trình đẳng áp}
\begin{ex}
% ID: L12C2B10-19-DS
% Mức: TH
Xét các phát biểu sau trong dạng “Tính thể tích và nhiệt độ trong quá trình đẳng áp”.
\choiceTF
{\True Khí có thể tích 7 lít ở 270 $K$, đun đẳng áp đến 340 K. Thể tích sau gần bằng — phát biểu đúng là: 8,81 lít.}
{Khí có thể tích 7 lít ở 270 $K$, đun đẳng áp đến 340 K. Thể tích sau gần bằng — phát biểu 5,56 lít là đúng.}
{\True Khí đẳng áp tăng nhiệt độ tuyệt đối từ 280 $K$ lên 360 K. Tỉ số $V_2/V_1$ bằng — phát biểu đúng là: 1,29.}
{Khí đẳng áp tăng nhiệt độ tuyệt đối từ 280 $K$ lên 360 K. Tỉ số $V_2/V_1$ bằng — phát biểu 1 là đúng.}
\loigiai{
\begin{itemchoice}
\item (Đúng) Khí có thể tích 7 lít ở 270 $K$, đun đẳng áp đến 340 K. Thể tích sau gần bằng — phát biểu đúng là: 8,81 lít. (Vì): $V_2=V_1T_2/T_1=7\cdot340/270=8,81$ lít. b) Sai: trái với hệ thức hoặc điều kiện đã nêu. c) Đúng: $V_2/V_1=T_2/T_1=1,29$. d) Sai: trái với hệ thức hoặc điều kiện đã nêu
\item (Sai) Khí có thể tích 7 lít ở 270 $K$, đun đẳng áp đến 340 K. Thể tích sau gần bằng — phát biểu 5,56 lít là đúng.
\item (Đúng) Khí đẳng áp tăng nhiệt độ tuyệt đối từ 280 $K$ lên 360 K. Tỉ số $V_2/V_1$ bằng — phát biểu đúng là: 1,29.
\item (Sai) Khí đẳng áp tăng nhiệt độ tuyệt đối từ 280 $K$ lên 360 K. Tỉ số $V_2/V_1$ bằng — phát biểu 1 là đúng.
\end{itemchoice}
}
\end{ex}

% ===== Câu 74 =====
% ID: L12C2B10-19-TL
% Mức: TH
\begin{ex}
% ID: L12C2B10-19-TL
% Mức: TH
Trong dạng “Tính thể tích và nhiệt độ trong quá trình đẳng áp”, Trình bày cơ sở lí thuyết của dạng “Tính thể tích và nhiệt độ trong quá trình đẳng áp” và nêu điều kiện áp dụng.
\choiceTF

\loigiai{
Nêu đúng đại lượng không đổi, hệ thức đặc trưng và yêu cầu dùng nhiệt độ kelvin, áp suất tuyệt đối khi có liên quan.
}
\end{ex}

% ===== Câu 75 =====
% ID: L12C2B10-19-TLN
% Mức: VD
\dangbt{Nhận biết quá trình đẳng áp và định luật Charles}
\begin{ex}
% ID: L12C2B10-19-TLN
% Mức: VD
Trong dạng “Nhận biết quá trình đẳng áp và định luật Charles”, Nhiệt độ tuyệt đối tăng 4 lần ở áp suất không đổi. Khối lượng riêng giảm bao nhiêu lần?
\shortans{4}
\loigiai{
$\rho\sim1/T$, nên khối lượng riêng giảm 4 lần.
}
\end{ex}

% ===== Câu 76 =====
% ID: L12C2B10-19-TN
% Mức: TH
\dangbt{Khối lượng riêng của chất khí trong quá trình đẳng áp}
\begin{ex}
% ID: L12C2B10-19-TN
% Mức: TH
Trong dạng “Khối lượng riêng của chất khí trong quá trình đẳng áp”, Đồ thị V- $T$ của quá trình đẳng áp trong thang kelvin là
\choice
{\True đường thẳng qua gốc}
{hypebol}
{đường ngang}
{đường thẳng không qua gốc}
\loigiai{
Đáp án A. $V$ tỉ lệ thuận $T$ tuyệt đối.
}
\end{ex}

% ===== Câu 77 =====
% ID: L12C2B10-20-DS
% Mức: VD
\dangbt{Tính thể tích và nhiệt độ trong quá trình đẳng áp}
\begin{ex}
% ID: L12C2B10-20-DS
% Mức: VD
Xét các phát biểu sau trong dạng “Tính thể tích và nhiệt độ trong quá trình đẳng áp”.
\choiceTF
{\True Đồ thị V- $T$ của quá trình đẳng áp trong thang kelvin là — phát biểu đúng là: đường thẳng qua gốc.}
{Đồ thị V- $T$ của quá trình đẳng áp trong thang kelvin là — phát biểu hypebol là đúng.}
{\True Khi áp suất không đổi và nhiệt độ tuyệt đối tăng gấp đôi, khối lượng riêng — phát biểu đúng là: giảm một nửa.}
{Khi áp suất không đổi và nhiệt độ tuyệt đối tăng gấp đôi, khối lượng riêng — phát biểu không đổi là đúng.}
\loigiai{
\begin{itemchoice}
\item (Đúng) Đồ thị V- $T$ của quá trình đẳng áp trong thang kelvin là — phát biểu đúng là: đường thẳng qua gốc. (Vì): $V$ tỉ lệ thuận $T$ tuyệt đối. b) Sai: trái với hệ thức hoặc điều kiện đã nêu. c) Đúng: $\rho=m/V$ và $V\sim T$, nên $\rho\sim1/T$. d) Sai: trái với hệ thức hoặc điều kiện đã nêu
\item (Sai) Đồ thị V- $T$ của quá trình đẳng áp trong thang kelvin là — phát biểu hypebol là đúng.
\item (Đúng) Khi áp suất không đổi và nhiệt độ tuyệt đối tăng gấp đôi, khối lượng riêng — phát biểu đúng là: giảm một nửa.
\item (Sai) Khi áp suất không đổi và nhiệt độ tuyệt đối tăng gấp đôi, khối lượng riêng — phát biểu không đổi là đúng.
\end{itemchoice}
}
\end{ex}

% ===== Câu 78 =====
% ID: L12C2B10-20-TL
% Mức: TH
\begin{ex}
% ID: L12C2B10-20-TL
% Mức: TH
Trong dạng “Tính thể tích và nhiệt độ trong quá trình đẳng áp”, Mô tả dạng đồ thị đặc trưng và cách nhận biết quá trình trong dạng “Tính thể tích và nhiệt độ trong quá trình đẳng áp”.
\choiceTF

\loigiai{
Xác định các trục, đại lượng không đổi và suy ra hình dạng đồ thị từ hệ thức vật lí tương ứng.
}
\end{ex}

% ===== Câu 79 =====
% ID: L12C2B10-20-TLN
% Mức: VDC
\dangbt{Nhận biết quá trình đẳng áp và định luật Charles}
\begin{ex}
% ID: L12C2B10-20-TLN
% Mức: VDC
Trong dạng “Nhận biết quá trình đẳng áp và định luật Charles”, Khí 16 lít ở 360 $K$ được làm lạnh đẳng áp về 300 K. Tính thể tích theo lít.
\shortans{13,33}
\loigiai{
$V_2=16\cdot300/360=13,33$ lít.
}
\end{ex}

% ===== Câu 80 =====
% ID: L12C2B10-20-TN
% Mức: VD
\dangbt{Khối lượng riêng của chất khí trong quá trình đẳng áp}
\begin{ex}
% ID: L12C2B10-20-TN
% Mức: VD
Trong dạng “Khối lượng riêng của chất khí trong quá trình đẳng áp”, Khi áp suất không đổi và nhiệt độ tuyệt đối tăng gấp đôi, khối lượng riêng
\choice
{\True giảm một nửa}
{tăng gấp đôi}
{không đổi}
{tăng bốn lần}
\loigiai{
Đáp án A. $\rho=m/V$ và $V\sim T$, nên $\rho\sim1/T$.
}
\end{ex}

% ===== Câu 81 =====
% ID: L12C2B10-21-DS
% Mức: VD
\dangbt{Tính thể tích và nhiệt độ trong quá trình đẳng áp}
\begin{ex}
% ID: L12C2B10-21-DS
% Mức: VD
Xét các phát biểu sau trong dạng “Tính thể tích và nhiệt độ trong quá trình đẳng áp”.
\choiceTF
{\True Trong định luật Charles, nhiệt độ phải đo theo — phát biểu đúng là: kelvin.}
{Trong định luật Charles, nhiệt độ phải đo theo — phát biểu độ $C$ tùy ý là đúng.}
{\True Quả bóng mềm căng hơn khi để nơi ấm, nếu áp suất gần không đổi, chủ yếu vì — phát biểu đúng là: thể tích khí tăng theo nhiệt độ.}
{Quả bóng mềm căng hơn khi để nơi ấm, nếu áp suất gần không đổi, chủ yếu vì — phát biểu khối lượng khí tăng là đúng.}
\loigiai{
\begin{itemchoice}
\item (Đúng) Trong định luật Charles, nhiệt độ phải đo theo — phát biểu đúng là: kelvin. (Vì): Tỉ lệ thuận chỉ đúng trực tiếp với nhiệt độ tuyệt đối. b) Sai: trái với hệ thức hoặc điều kiện đã nêu. c) Đúng: Theo Charles, khi $T$ tăng thì $V$ tăng ở áp suất không đổi. d) Sai: trái với hệ thức hoặc điều kiện đã nêu
\item (Sai) Trong định luật Charles, nhiệt độ phải đo theo — phát biểu độ $C$ tùy ý là đúng.
\item (Đúng) Quả bóng mềm căng hơn khi để nơi ấm, nếu áp suất gần không đổi, chủ yếu vì — phát biểu đúng là: thể tích khí tăng theo nhiệt độ.
\item (Sai) Quả bóng mềm căng hơn khi để nơi ấm, nếu áp suất gần không đổi, chủ yếu vì — phát biểu khối lượng khí tăng là đúng.
\end{itemchoice}
}
\end{ex}

% ===== Câu 82 =====
% ID: L12C2B10-21-TL
% Mức: VD
\begin{bt}
% ID: L12C2B10-21-TL
% Mức: VD
Trong dạng “Tính thể tích và nhiệt độ trong quá trình đẳng áp”, Nhiệt độ tuyệt đối tăng từ 280 $K$ lên 340 $K$ ở áp suất không đổi. Tính $V_2/V_1$.
\loigiai{
$V_2/V_1=T_2/T_1=1,21$.
}
\end{bt}

% ===== Câu 83 =====
% ID: L12C2B10-21-TLN
% Mức: TH
\begin{ex}
% ID: L12C2B10-21-TLN
% Mức: TH
Trong dạng “Tính thể tích và nhiệt độ trong quá trình đẳng áp”, Khí 5 lít ở 310 $K$ được đun đẳng áp đến 370 K. Tính thể tích sau theo lít.
\shortans{5,97}
\loigiai{
$V_2=5\cdot370/310=5,97$ lít.
}
\end{ex}

% ===== Câu 84 =====
% ID: L12C2B10-21-TN
% Mức: VD
\dangbt{Khối lượng riêng của chất khí trong quá trình đẳng áp}
\begin{ex}
% ID: L12C2B10-21-TN
% Mức: VD
Trong dạng “Khối lượng riêng của chất khí trong quá trình đẳng áp”, Trong định luật Charles, nhiệt độ phải đo theo
\choice
{\True kelvin}
{độ $C$ tùy ý}
{độ $F$ tùy ý}
{mọi thang đều được}
\loigiai{
Đáp án A. Tỉ lệ thuận chỉ đúng trực tiếp với nhiệt độ tuyệt đối.
}
\end{ex}

% ===== Câu 85 =====
% ID: L12C2B10-22-DS
% Mức: VDC
\dangbt{Tính thể tích và nhiệt độ trong quá trình đẳng áp}
\begin{ex}
% ID: L12C2B10-22-DS
% Mức: VDC
Xét các phát biểu sau trong dạng “Tính thể tích và nhiệt độ trong quá trình đẳng áp”.
\choiceTF
{\True Đường đẳng áp có áp suất nhỏ hơn trên đồ thị V- $T$ của cùng lượng khí nằm — phát biểu đúng là: dốc hơn.}
{Đường đẳng áp có áp suất nhỏ hơn trên đồ thị V- $T$ của cùng lượng khí nằm — phát biểu thoải hơn là đúng.}
{\True Ở áp suất không đổi, thể tích ngoại suy bằng không tại — phát biểu đúng là: 0 K.}
{Ở áp suất không đổi, thể tích ngoại suy bằng không tại — phát biểu 100 ° $C$ là đúng.}
\loigiai{
\begin{itemchoice}
\item (Đúng) Đường đẳng áp có áp suất nhỏ hơn trên đồ thị V- $T$ của cùng lượng khí nằm — phát biểu đúng là: dốc hơn. (Vì): Từ $V=nRT/p$, p nhỏ hơn cho hệ số góc lớn hơn. b) Sai: trái với hệ thức hoặc điều kiện đã nêu. c) Đúng: 0 $K$ là độ không tuyệt đối theo mô hình khí lí tưởng. d) Sai: trái với hệ thức hoặc điều kiện đã nêu
\item (Sai) Đường đẳng áp có áp suất nhỏ hơn trên đồ thị V- $T$ của cùng lượng khí nằm — phát biểu thoải hơn là đúng.
\item (Đúng) Ở áp suất không đổi, thể tích ngoại suy bằng không tại — phát biểu đúng là: 0 K.
\item (Sai) Ở áp suất không đổi, thể tích ngoại suy bằng không tại — phát biểu 100 ° $C$ là đúng.
\end{itemchoice}
}
\end{ex}

% ===== Câu 86 =====
% ID: L12C2B10-22-TL
% Mức: VD
\begin{ex}
% ID: L12C2B10-22-TL
% Mức: VD
Trong dạng “Tính thể tích và nhiệt độ trong quá trình đẳng áp”, 4
\choiceTF

\loigiai{
Do $V/T$ không đổi nên $T_2/T_1=V_2/V_1=4$.
}
\end{ex}

% ===== Câu 87 =====
% ID: L12C2B10-22-TLN
% Mức: TH
\begin{ex}
% ID: L12C2B10-22-TLN
% Mức: TH
Trong dạng “Tính thể tích và nhiệt độ trong quá trình đẳng áp”, Khí có thể tích 7 lít ở 320 $K$, đẳng áp tăng thể tích lên 14 lít. Tính nhiệt độ sau theo K.
\shortans{640}
\loigiai{
$T_2=T_1V_2/V_1=320\cdot2=640$ K.
}
\end{ex}

% ===== Câu 88 =====
% ID: L12C2B10-22-TN
% Mức: VD
\dangbt{Khối lượng riêng của chất khí trong quá trình đẳng áp}
\begin{ex}
% ID: L12C2B10-22-TN
% Mức: VD
Trong dạng “Khối lượng riêng của chất khí trong quá trình đẳng áp”, Quả bóng mềm căng hơn khi để nơi ấm, nếu áp suất gần không đổi, chủ yếu vì
\choice
{\True thể tích khí tăng theo nhiệt độ}
{số mol khí tăng}
{khối lượng khí tăng}
{áp suất bằng không}
\loigiai{
Đáp án A. Theo Charles, khi $T$ tăng thì $V$ tăng ở áp suất không đổi.
}
\end{ex}

% ===== Câu 89 =====
% ID: L12C2B10-23-DS
% Mức: TH
\dangbt{Đồ thị đẳng áp trên các hệ tọa độ}
\begin{ex}
% ID: L12C2B10-23-DS
% Mức: TH
Xét các phát biểu sau trong dạng “Đồ thị đẳng áp trên các hệ tọa độ”.
\choiceTF
{\True Trong quá trình đẳng áp của một lượng khí xác định, đại lượng không đổi là — phát biểu đúng là: áp suất.}
{Trong quá trình đẳng áp của một lượng khí xác định, đại lượng không đổi là — phát biểu thể tích là đúng.}
{\True Hệ thức đúng của định luật Charles là — phát biểu đúng là: $V_1/T_1=V_2/T_2$.}
{Hệ thức đúng của định luật Charles là — phát biểu $p_1/T_1=p_2/T_2$ là đúng.}
\loigiai{
\begin{itemchoice}
\item (Đúng) Trong quá trình đẳng áp của một lượng khí xác định, đại lượng không đổi là — phát biểu đúng là: áp suất. (Vì): Quá trình đẳng áp có áp suất không đổi. b) Sai: trái với hệ thức hoặc điều kiện đã nêu. c) Đúng: Ở áp suất không đổi, $V/T$ là hằng số. d) Sai: trái với hệ thức hoặc điều kiện đã nêu
\item (Sai) Trong quá trình đẳng áp của một lượng khí xác định, đại lượng không đổi là — phát biểu thể tích là đúng.
\item (Đúng) Hệ thức đúng của định luật Charles là — phát biểu đúng là: $V_1/T_1=V_2/T_2$.
\item (Sai) Hệ thức đúng của định luật Charles là — phát biểu $p_1/T_1=p_2/T_2$ là đúng.
\end{itemchoice}
}
\end{ex}

% ===== Câu 90 =====
% ID: L12C2B10-23-TL
% Mức: VD
\dangbt{Tính thể tích và nhiệt độ trong quá trình đẳng áp}
\begin{bt}
% ID: L12C2B10-23-TL
% Mức: VD
Trong dạng “Tính thể tích và nhiệt độ trong quá trình đẳng áp”, $\rho\sim1/T$, nên khối lượng riêng giảm 2 lần.
\loigiai{
$\rho\sim1/T$, nên khối lượng riêng giảm 2 lần.
}
\end{bt}

% ===== Câu 91 =====
% ID: L12C2B10-23-TLN
% Mức: VD
\begin{ex}
% ID: L12C2B10-23-TLN
% Mức: VD
Trong dạng “Tính thể tích và nhiệt độ trong quá trình đẳng áp”, Nhiệt độ tuyệt đối tăng từ 280 $K$ lên 340 $K$ ở áp suất không đổi. Tính $V_2/V_1$.
\shortans{1,21}
\loigiai{
$V_2/V_1=T_2/T_1=1,21$.
}
\end{ex}

% ===== Câu 92 =====
% ID: L12C2B10-23-TN
% Mức: VDC
\dangbt{Khối lượng riêng của chất khí trong quá trình đẳng áp}
\begin{ex}
% ID: L12C2B10-23-TN
% Mức: VDC
Trong dạng “Khối lượng riêng của chất khí trong quá trình đẳng áp”, Đường đẳng áp có áp suất nhỏ hơn trên đồ thị V- $T$ của cùng lượng khí nằm
\choice
{\True dốc hơn}
{thoải hơn}
{trùng trục $T$}
{luôn trùng nhau}
\loigiai{
Đáp án A. Từ $V=nRT/p$, p nhỏ hơn cho hệ số góc lớn hơn.
}
\end{ex}

% ===== Câu 93 =====
% ID: L12C2B10-24-DS
% Mức: TH
\dangbt{Đồ thị đẳng áp trên các hệ tọa độ}
\begin{ex}
% ID: L12C2B10-24-DS
% Mức: TH
Xét các phát biểu sau trong dạng “Đồ thị đẳng áp trên các hệ tọa độ”.
\choiceTF
{\True Khí có thể tích 8 lít ở 280 $K$, đun đẳng áp đến 350 K. Thể tích sau gần bằng — phát biểu đúng là: 10 lít.}
{Khí có thể tích 8 lít ở 280 $K$, đun đẳng áp đến 350 K. Thể tích sau gần bằng — phát biểu 6,4 lít là đúng.}
{\True Khí đẳng áp tăng nhiệt độ tuyệt đối từ 290 $K$ lên 370 K. Tỉ số $V_2/V_1$ bằng — phát biểu đúng là: 1,28.}
{Khí đẳng áp tăng nhiệt độ tuyệt đối từ 290 $K$ lên 370 K. Tỉ số $V_2/V_1$ bằng — phát biểu 1 là đúng.}
\loigiai{
\begin{itemchoice}
\item (Đúng) Khí có thể tích 8 lít ở 280 $K$, đun đẳng áp đến 350 K. Thể tích sau gần bằng — phát biểu đúng là: 10 lít. (Vì): $V_2=V_1T_2/T_1=8\cdot350/280=10$ lít. b) Sai: trái với hệ thức hoặc điều kiện đã nêu. c) Đúng: $V_2/V_1=T_2/T_1=1,28$. d) Sai: trái với hệ thức hoặc điều kiện đã nêu
\item (Sai) Khí có thể tích 8 lít ở 280 $K$, đun đẳng áp đến 350 K. Thể tích sau gần bằng — phát biểu 6,4 lít là đúng.
\item (Đúng) Khí đẳng áp tăng nhiệt độ tuyệt đối từ 290 $K$ lên 370 K. Tỉ số $V_2/V_1$ bằng — phát biểu đúng là: 1,28.
\item (Sai) Khí đẳng áp tăng nhiệt độ tuyệt đối từ 290 $K$ lên 370 K. Tỉ số $V_2/V_1$ bằng — phát biểu 1 là đúng.
\end{itemchoice}
}
\end{ex}

% ===== Câu 94 =====
% ID: L12C2B10-24-TL
% Mức: VDC
\dangbt{Tính thể tích và nhiệt độ trong quá trình đẳng áp}
\begin{bt}
% ID: L12C2B10-24-TL
% Mức: VDC
Trong dạng “Tính thể tích và nhiệt độ trong quá trình đẳng áp”, undefined
\loigiai{
$V_2=9\cdot310/370=7,54$ lít.
}
\end{bt}

% ===== Câu 95 =====
% ID: L12C2B10-24-TLN
% Mức: VD
\begin{ex}
% ID: L12C2B10-24-TLN
% Mức: VD
Trong dạng “Tính thể tích và nhiệt độ trong quá trình đẳng áp”, Thể tích đẳng áp tăng 4 lần. Nhiệt độ tuyệt đối tăng bao nhiêu lần?
\shortans{4}
\loigiai{
Do $V/T$ không đổi nên $T_2/T_1=V_2/V_1=4$.
}
\end{ex}

% ===== Câu 96 =====
% ID: L12C2B10-24-TN
% Mức: VDC
\dangbt{Khối lượng riêng của chất khí trong quá trình đẳng áp}
\begin{ex}
% ID: L12C2B10-24-TN
% Mức: VDC
Trong dạng “Khối lượng riêng của chất khí trong quá trình đẳng áp”, Ở áp suất không đổi, thể tích ngoại suy bằng không tại
\choice
{\True 0 $K$}
{0 ° $C$}
{100 ° $C$}
{273 $K$}
\loigiai{
Đáp án A. 0 $K$ là độ không tuyệt đối theo mô hình khí lí tưởng.
}
\end{ex}

% ===== Câu 97 =====
% ID: L12C2B10-25-DS
% Mức: VD
\dangbt{Đồ thị đẳng áp trên các hệ tọa độ}
\begin{ex}
% ID: L12C2B10-25-DS
% Mức: VD
Xét các phát biểu sau trong dạng “Đồ thị đẳng áp trên các hệ tọa độ”.
\choiceTF
{\True Đồ thị V- $T$ của quá trình đẳng áp trong thang kelvin là — phát biểu đúng là: đường thẳng qua gốc.}
{Đồ thị V- $T$ của quá trình đẳng áp trong thang kelvin là — phát biểu hypebol là đúng.}
{\True Khi áp suất không đổi và nhiệt độ tuyệt đối tăng gấp đôi, khối lượng riêng — phát biểu đúng là: giảm một nửa.}
{Khi áp suất không đổi và nhiệt độ tuyệt đối tăng gấp đôi, khối lượng riêng — phát biểu không đổi là đúng.}
\loigiai{
\begin{itemchoice}
\item (Đúng) Đồ thị V- $T$ của quá trình đẳng áp trong thang kelvin là — phát biểu đúng là: đường thẳng qua gốc. (Vì): $V$ tỉ lệ thuận $T$ tuyệt đối. b) Sai: trái với hệ thức hoặc điều kiện đã nêu. c) Đúng: $\rho=m/V$ và $V\sim T$, nên $\rho\sim1/T$. d) Sai: trái với hệ thức hoặc điều kiện đã nêu
\item (Sai) Đồ thị V- $T$ của quá trình đẳng áp trong thang kelvin là — phát biểu hypebol là đúng.
\item (Đúng) Khi áp suất không đổi và nhiệt độ tuyệt đối tăng gấp đôi, khối lượng riêng — phát biểu đúng là: giảm một nửa.
\item (Sai) Khi áp suất không đổi và nhiệt độ tuyệt đối tăng gấp đôi, khối lượng riêng — phát biểu không đổi là đúng.
\end{itemchoice}
}
\end{ex}

% ===== Câu 98 =====
% ID: L12C2B10-25-TL
% Mức: TH
\begin{ex}
% ID: L12C2B10-25-TL
% Mức: TH
Trong dạng “Đồ thị đẳng áp trên các hệ tọa độ”, Trình bày cơ sở lí thuyết của dạng “Đồ thị đẳng áp trên các hệ tọa độ” và nêu điều kiện áp dụng.
\choiceTF

\loigiai{
Nêu đúng đại lượng không đổi, hệ thức đặc trưng và yêu cầu dùng nhiệt độ kelvin, áp suất tuyệt đối khi có liên quan.
}
\end{ex}

% ===== Câu 99 =====
% ID: L12C2B10-25-TLN
% Mức: VD
\dangbt{Tính thể tích và nhiệt độ trong quá trình đẳng áp}
\begin{ex}
% ID: L12C2B10-25-TLN
% Mức: VD
Trong dạng “Tính thể tích và nhiệt độ trong quá trình đẳng áp”, Nhiệt độ tuyệt đối tăng 2 lần ở áp suất không đổi. Khối lượng riêng giảm bao nhiêu lần?
\shortans{2}
\loigiai{
$\rho\sim1/T$, nên khối lượng riêng giảm 2 lần.
}
\end{ex}

% ===== Câu 100 =====
% ID: L12C2B10-25-TN
% Mức: NB
\dangbt{Nhận biết quá trình đẳng áp và định luật Charles}
\begin{ex}
% ID: L12C2B10-25-TN
% Mức: NB
Trong dạng “Nhận biết quá trình đẳng áp và định luật Charles”, Trong quá trình đẳng áp của một lượng khí xác định, đại lượng không đổi là
\choice
{\True áp suất}
{thể tích}
{nhiệt độ}
{khối lượng riêng}
\loigiai{
Đáp án A. Quá trình đẳng áp có áp suất không đổi.
}
\end{ex}

% ===== Câu 101 =====
% ID: L12C2B10-26-DS
% Mức: VD
\dangbt{Đồ thị đẳng áp trên các hệ tọa độ}
\begin{ex}
% ID: L12C2B10-26-DS
% Mức: VD
Xét các phát biểu sau trong dạng “Đồ thị đẳng áp trên các hệ tọa độ”.
\choiceTF
{\True Trong định luật Charles, nhiệt độ phải đo theo — phát biểu đúng là: kelvin.}
{Trong định luật Charles, nhiệt độ phải đo theo — phát biểu độ $C$ tùy ý là đúng.}
{\True Quả bóng mềm căng hơn khi để nơi ấm, nếu áp suất gần không đổi, chủ yếu vì — phát biểu đúng là: thể tích khí tăng theo nhiệt độ.}
{Quả bóng mềm căng hơn khi để nơi ấm, nếu áp suất gần không đổi, chủ yếu vì — phát biểu khối lượng khí tăng là đúng.}
\loigiai{
\begin{itemchoice}
\item (Đúng) Trong định luật Charles, nhiệt độ phải đo theo — phát biểu đúng là: kelvin. (Vì): Tỉ lệ thuận chỉ đúng trực tiếp với nhiệt độ tuyệt đối. b) Sai: trái với hệ thức hoặc điều kiện đã nêu. c) Đúng: Theo Charles, khi $T$ tăng thì $V$ tăng ở áp suất không đổi. d) Sai: trái với hệ thức hoặc điều kiện đã nêu
\item (Sai) Trong định luật Charles, nhiệt độ phải đo theo — phát biểu độ $C$ tùy ý là đúng.
\item (Đúng) Quả bóng mềm căng hơn khi để nơi ấm, nếu áp suất gần không đổi, chủ yếu vì — phát biểu đúng là: thể tích khí tăng theo nhiệt độ.
\item (Sai) Quả bóng mềm căng hơn khi để nơi ấm, nếu áp suất gần không đổi, chủ yếu vì — phát biểu khối lượng khí tăng là đúng.
\end{itemchoice}
}
\end{ex}

% ===== Câu 102 =====
% ID: L12C2B10-26-TL
% Mức: TH
\begin{ex}
% ID: L12C2B10-26-TL
% Mức: TH
Trong dạng “Đồ thị đẳng áp trên các hệ tọa độ”, Mô tả dạng đồ thị đặc trưng và cách nhận biết quá trình trong dạng “Đồ thị đẳng áp trên các hệ tọa độ”.
\choiceTF

\loigiai{
Xác định các trục, đại lượng không đổi và suy ra hình dạng đồ thị từ hệ thức vật lí tương ứng.
}
\end{ex}

% ===== Câu 103 =====
% ID: L12C2B10-26-TLN
% Mức: VDC
\dangbt{Tính thể tích và nhiệt độ trong quá trình đẳng áp}
\begin{ex}
% ID: L12C2B10-26-TLN
% Mức: VDC
Trong dạng “Tính thể tích và nhiệt độ trong quá trình đẳng áp”, Khí 9 lít ở 370 $K$ được làm lạnh đẳng áp về 310 K. Tính thể tích theo lít.
\shortans{7,54}
\loigiai{
$V_2=9\cdot310/370=7,54$ lít.
}
\end{ex}

% ===== Câu 104 =====
% ID: L12C2B10-26-TN
% Mức: NB
\dangbt{Nhận biết quá trình đẳng áp và định luật Charles}
\begin{ex}
% ID: L12C2B10-26-TN
% Mức: NB
Trong dạng “Nhận biết quá trình đẳng áp và định luật Charles”, Hệ thức đúng của định luật Charles là
\choice
{\True $V_1/T_1=V_2/T_2$}
{$p_1V_1=p_2V_2$}
{$p_1/T_1=p_2/T_2$}
{$V_1T_1=V_2T_2$}
\loigiai{
Đáp án A. Ở áp suất không đổi, $V/T$ là hằng số.
}
\end{ex}

% ===== Câu 105 =====
% ID: L12C2B10-27-DS
% Mức: VDC
\dangbt{Đồ thị đẳng áp trên các hệ tọa độ}
\begin{ex}
% ID: L12C2B10-27-DS
% Mức: VDC
Xét các phát biểu sau trong dạng “Đồ thị đẳng áp trên các hệ tọa độ”.
\choiceTF
{\True Đường đẳng áp có áp suất nhỏ hơn trên đồ thị V- $T$ của cùng lượng khí nằm — phát biểu đúng là: dốc hơn.}
{Đường đẳng áp có áp suất nhỏ hơn trên đồ thị V- $T$ của cùng lượng khí nằm — phát biểu thoải hơn là đúng.}
{\True Ở áp suất không đổi, thể tích ngoại suy bằng không tại — phát biểu đúng là: 0 K.}
{Ở áp suất không đổi, thể tích ngoại suy bằng không tại — phát biểu 100 ° $C$ là đúng.}
\loigiai{
\begin{itemchoice}
\item (Đúng) Đường đẳng áp có áp suất nhỏ hơn trên đồ thị V- $T$ của cùng lượng khí nằm — phát biểu đúng là: dốc hơn. (Vì): Từ $V=nRT/p$, p nhỏ hơn cho hệ số góc lớn hơn. b) Sai: trái với hệ thức hoặc điều kiện đã nêu. c) Đúng: 0 $K$ là độ không tuyệt đối theo mô hình khí lí tưởng. d) Sai: trái với hệ thức hoặc điều kiện đã nêu
\item (Sai) Đường đẳng áp có áp suất nhỏ hơn trên đồ thị V- $T$ của cùng lượng khí nằm — phát biểu thoải hơn là đúng.
\item (Đúng) Ở áp suất không đổi, thể tích ngoại suy bằng không tại — phát biểu đúng là: 0 K.
\item (Sai) Ở áp suất không đổi, thể tích ngoại suy bằng không tại — phát biểu 100 ° $C$ là đúng.
\end{itemchoice}
}
\end{ex}

% ===== Câu 106 =====
% ID: L12C2B10-27-TL
% Mức: VD
\begin{bt}
% ID: L12C2B10-27-TL
% Mức: VD
Trong dạng “Đồ thị đẳng áp trên các hệ tọa độ”, Nhiệt độ tuyệt đối tăng từ 290 $K$ lên 350 $K$ ở áp suất không đổi. Tính $V_2/V_1$.
\loigiai{
$V_2/V_1=T_2/T_1=1,21$.
}
\end{bt}

% ===== Câu 107 =====
% ID: L12C2B10-27-TLN
% Mức: TH
\begin{ex}
% ID: L12C2B10-27-TLN
% Mức: TH
Trong dạng “Đồ thị đẳng áp trên các hệ tọa độ”, Khí 6 lít ở 320 $K$ được đun đẳng áp đến 380 K. Tính thể tích sau theo lít.
\shortans{7,13}
\loigiai{
$V_2=6\cdot380/320=7,13$ lít.
}
\end{ex}

% ===== Câu 108 =====
% ID: L12C2B10-27-TN
% Mức: TH
\dangbt{Nhận biết quá trình đẳng áp và định luật Charles}
\begin{ex}
% ID: L12C2B10-27-TN
% Mức: TH
Trong dạng “Nhận biết quá trình đẳng áp và định luật Charles”, Khí có thể tích 6 lít ở 260 $K$, đun đẳng áp đến 330 K. Thể tích sau gần bằng
\choice
{\True 7,62 lít}
{4,73 lít}
{6 lít}
{15,23 lít}
\loigiai{
Đáp án A. $V_2=V_1T_2/T_1=6\cdot330/260=7,62$ lít.
}
\end{ex}

% ===== Câu 109 =====
% ID: L12C2B10-28-DS
% Mức: TH
\dangbt{Ứng dụng định luật Charles trong thực tế}
\begin{ex}
% ID: L12C2B10-28-DS
% Mức: TH
Xét các phát biểu sau trong dạng “Ứng dụng định luật Charles trong thực tế”.
\choiceTF
{\True Trong quá trình đẳng áp của một lượng khí xác định, đại lượng không đổi là — phát biểu đúng là: áp suất.}
{Trong quá trình đẳng áp của một lượng khí xác định, đại lượng không đổi là — phát biểu thể tích là đúng.}
{\True Hệ thức đúng của định luật Charles là — phát biểu đúng là: $V_1/T_1=V_2/T_2$.}
{Hệ thức đúng của định luật Charles là — phát biểu $p_1/T_1=p_2/T_2$ là đúng.}
\loigiai{
\begin{itemchoice}
\item (Đúng) Trong quá trình đẳng áp của một lượng khí xác định, đại lượng không đổi là — phát biểu đúng là: áp suất. (Vì): Quá trình đẳng áp có áp suất không đổi. b) Sai: trái với hệ thức hoặc điều kiện đã nêu. c) Đúng: Ở áp suất không đổi, $V/T$ là hằng số. d) Sai: trái với hệ thức hoặc điều kiện đã nêu
\item (Sai) Trong quá trình đẳng áp của một lượng khí xác định, đại lượng không đổi là — phát biểu thể tích là đúng.
\item (Đúng) Hệ thức đúng của định luật Charles là — phát biểu đúng là: $V_1/T_1=V_2/T_2$.
\item (Sai) Hệ thức đúng của định luật Charles là — phát biểu $p_1/T_1=p_2/T_2$ là đúng.
\end{itemchoice}
}
\end{ex}

% ===== Câu 110 =====
% ID: L12C2B10-28-TL
% Mức: VD
\dangbt{Đồ thị đẳng áp trên các hệ tọa độ}
\begin{ex}
% ID: L12C2B10-28-TL
% Mức: VD
Trong dạng “Đồ thị đẳng áp trên các hệ tọa độ”, 2
\choiceTF

\loigiai{
Do $V/T$ không đổi nên $T_2/T_1=V_2/V_1=2$.
}
\end{ex}

% ===== Câu 111 =====
% ID: L12C2B10-28-TLN
% Mức: TH
\begin{ex}
% ID: L12C2B10-28-TLN
% Mức: TH
Trong dạng “Đồ thị đẳng áp trên các hệ tọa độ”, Khí có thể tích 8 lít ở 280 $K$, đẳng áp tăng thể tích lên 24 lít. Tính nhiệt độ sau theo K.
\shortans{840}
\loigiai{
$T_2=T_1V_2/V_1=280\cdot3=840$ K.
}
\end{ex}

% ===== Câu 112 =====
% ID: L12C2B10-28-TN
% Mức: TH
\dangbt{Nhận biết quá trình đẳng áp và định luật Charles}
\begin{ex}
% ID: L12C2B10-28-TN
% Mức: TH
Trong dạng “Nhận biết quá trình đẳng áp và định luật Charles”, Khí đẳng áp tăng nhiệt độ tuyệt đối từ 270 $K$ lên 350 K. Tỉ số $V_2/V_1$ bằng
\choice
{\True 1,3}
{0,77}
{1}
{0,3}
\loigiai{
Đáp án A. $V_2/V_1=T_2/T_1=1,3$.
}
\end{ex}

% ===== Câu 113 =====
% ID: L12C2B10-29-DS
% Mức: TH
\dangbt{Ứng dụng định luật Charles trong thực tế}
\begin{ex}
% ID: L12C2B10-29-DS
% Mức: TH
Xét các phát biểu sau trong dạng “Ứng dụng định luật Charles trong thực tế”.
\choiceTF
{\True Khí có thể tích 3 lít ở 300 $K$, đun đẳng áp đến 370 K. Thể tích sau gần bằng — phát biểu đúng là: 3,7 lít.}
{Khí có thể tích 3 lít ở 300 $K$, đun đẳng áp đến 370 K. Thể tích sau gần bằng — phát biểu 2,43 lít là đúng.}
{\True Khí đẳng áp tăng nhiệt độ tuyệt đối từ 310 $K$ lên 390 K. Tỉ số $V_2/V_1$ bằng — phát biểu đúng là: 1,26.}
{Khí đẳng áp tăng nhiệt độ tuyệt đối từ 310 $K$ lên 390 K. Tỉ số $V_2/V_1$ bằng — phát biểu 1 là đúng.}
\loigiai{
\begin{itemchoice}
\item (Đúng) Khí có thể tích 3 lít ở 300 $K$, đun đẳng áp đến 370 K. Thể tích sau gần bằng — phát biểu đúng là: 3,7 lít. (Vì): $V_2=V_1T_2/T_1=3\cdot370/300=3,7$ lít. b) Sai: trái với hệ thức hoặc điều kiện đã nêu. c) Đúng: $V_2/V_1=T_2/T_1=1,26$. d) Sai: trái với hệ thức hoặc điều kiện đã nêu
\item (Sai) Khí có thể tích 3 lít ở 300 $K$, đun đẳng áp đến 370 K. Thể tích sau gần bằng — phát biểu 2,43 lít là đúng.
\item (Đúng) Khí đẳng áp tăng nhiệt độ tuyệt đối từ 310 $K$ lên 390 K. Tỉ số $V_2/V_1$ bằng — phát biểu đúng là: 1,26.
\item (Sai) Khí đẳng áp tăng nhiệt độ tuyệt đối từ 310 $K$ lên 390 K. Tỉ số $V_2/V_1$ bằng — phát biểu 1 là đúng.
\end{itemchoice}
}
\end{ex}

% ===== Câu 114 =====
% ID: L12C2B10-29-TL
% Mức: VD
\dangbt{Đồ thị đẳng áp trên các hệ tọa độ}
\begin{bt}
% ID: L12C2B10-29-TL
% Mức: VD
Trong dạng “Đồ thị đẳng áp trên các hệ tọa độ”, $\rho\sim1/T$, nên khối lượng riêng giảm 3 lần.
\loigiai{
$\rho\sim1/T$, nên khối lượng riêng giảm 3 lần.
}
\end{bt}

% ===== Câu 115 =====
% ID: L12C2B10-29-TLN
% Mức: VD
\begin{ex}
% ID: L12C2B10-29-TLN
% Mức: VD
Trong dạng “Đồ thị đẳng áp trên các hệ tọa độ”, Nhiệt độ tuyệt đối tăng từ 290 $K$ lên 350 $K$ ở áp suất không đổi. Tính $V_2/V_1$.
\shortans{1,21}
\loigiai{
$V_2/V_1=T_2/T_1=1,21$.
}
\end{ex}

% ===== Câu 116 =====
% ID: L12C2B10-29-TN
% Mức: TH
\dangbt{Nhận biết quá trình đẳng áp và định luật Charles}
\begin{ex}
% ID: L12C2B10-29-TN
% Mức: TH
Trong dạng “Nhận biết quá trình đẳng áp và định luật Charles”, Đồ thị V- $T$ của quá trình đẳng áp trong thang kelvin là
\choice
{\True đường thẳng qua gốc}
{hypebol}
{đường ngang}
{đường thẳng không qua gốc}
\loigiai{
Đáp án A. $V$ tỉ lệ thuận $T$ tuyệt đối.
}
\end{ex}

% ===== Câu 117 =====
% ID: L12C2B10-30-DS
% Mức: VD
\dangbt{Ứng dụng định luật Charles trong thực tế}
\begin{ex}
% ID: L12C2B10-30-DS
% Mức: VD
Xét các phát biểu sau trong dạng “Ứng dụng định luật Charles trong thực tế”.
\choiceTF
{\True Đồ thị V- $T$ của quá trình đẳng áp trong thang kelvin là — phát biểu đúng là: đường thẳng qua gốc.}
{Đồ thị V- $T$ của quá trình đẳng áp trong thang kelvin là — phát biểu hypebol là đúng.}
{\True Khi áp suất không đổi và nhiệt độ tuyệt đối tăng gấp đôi, khối lượng riêng — phát biểu đúng là: giảm một nửa.}
{Khi áp suất không đổi và nhiệt độ tuyệt đối tăng gấp đôi, khối lượng riêng — phát biểu không đổi là đúng.}
\loigiai{
\begin{itemchoice}
\item (Đúng) Đồ thị V- $T$ của quá trình đẳng áp trong thang kelvin là — phát biểu đúng là: đường thẳng qua gốc. (Vì): $V$ tỉ lệ thuận $T$ tuyệt đối. b) Sai: trái với hệ thức hoặc điều kiện đã nêu. c) Đúng: $\rho=m/V$ và $V\sim T$, nên $\rho\sim1/T$. d) Sai: trái với hệ thức hoặc điều kiện đã nêu
\item (Sai) Đồ thị V- $T$ của quá trình đẳng áp trong thang kelvin là — phát biểu hypebol là đúng.
\item (Đúng) Khi áp suất không đổi và nhiệt độ tuyệt đối tăng gấp đôi, khối lượng riêng — phát biểu đúng là: giảm một nửa.
\item (Sai) Khi áp suất không đổi và nhiệt độ tuyệt đối tăng gấp đôi, khối lượng riêng — phát biểu không đổi là đúng.
\end{itemchoice}
}
\end{ex}

% ===== Câu 118 =====
% ID: L12C2B10-30-TL
% Mức: VDC
\dangbt{Đồ thị đẳng áp trên các hệ tọa độ}
\begin{bt}
% ID: L12C2B10-30-TL
% Mức: VDC
Trong dạng “Đồ thị đẳng áp trên các hệ tọa độ”, undefined
\loigiai{
$V_2=16\cdot320/380=13,47$ lít.
}
\end{bt}

% ===== Câu 119 =====
% ID: L12C2B10-30-TLN
% Mức: VD
\begin{ex}
% ID: L12C2B10-30-TLN
% Mức: VD
Trong dạng “Đồ thị đẳng áp trên các hệ tọa độ”, Thể tích đẳng áp tăng 2 lần. Nhiệt độ tuyệt đối tăng bao nhiêu lần?
\shortans{2}
\loigiai{
Do $V/T$ không đổi nên $T_2/T_1=V_2/V_1=2$.
}
\end{ex}

% ===== Câu 120 =====
% ID: L12C2B10-30-TN
% Mức: VD
\dangbt{Nhận biết quá trình đẳng áp và định luật Charles}
\begin{ex}
% ID: L12C2B10-30-TN
% Mức: VD
Trong dạng “Nhận biết quá trình đẳng áp và định luật Charles”, Khi áp suất không đổi và nhiệt độ tuyệt đối tăng gấp đôi, khối lượng riêng
\choice
{\True giảm một nửa}
{tăng gấp đôi}
{không đổi}
{tăng bốn lần}
\loigiai{
Đáp án A. $\rho=m/V$ và $V\sim T$, nên $\rho\sim1/T$.
}
\end{ex}

% ===== Câu 121 =====
% ID: L12C2B10-31-DS
% Mức: VD
\dangbt{Ứng dụng định luật Charles trong thực tế}
\begin{ex}
% ID: L12C2B10-31-DS
% Mức: VD
Xét các phát biểu sau trong dạng “Ứng dụng định luật Charles trong thực tế”.
\choiceTF
{\True Trong định luật Charles, nhiệt độ phải đo theo — phát biểu đúng là: kelvin.}
{Trong định luật Charles, nhiệt độ phải đo theo — phát biểu độ $C$ tùy ý là đúng.}
{\True Quả bóng mềm căng hơn khi để nơi ấm, nếu áp suất gần không đổi, chủ yếu vì — phát biểu đúng là: thể tích khí tăng theo nhiệt độ.}
{Quả bóng mềm căng hơn khi để nơi ấm, nếu áp suất gần không đổi, chủ yếu vì — phát biểu khối lượng khí tăng là đúng.}
\loigiai{
\begin{itemchoice}
\item (Đúng) Trong định luật Charles, nhiệt độ phải đo theo — phát biểu đúng là: kelvin. (Vì): Tỉ lệ thuận chỉ đúng trực tiếp với nhiệt độ tuyệt đối. b) Sai: trái với hệ thức hoặc điều kiện đã nêu. c) Đúng: Theo Charles, khi $T$ tăng thì $V$ tăng ở áp suất không đổi. d) Sai: trái với hệ thức hoặc điều kiện đã nêu
\item (Sai) Trong định luật Charles, nhiệt độ phải đo theo — phát biểu độ $C$ tùy ý là đúng.
\item (Đúng) Quả bóng mềm căng hơn khi để nơi ấm, nếu áp suất gần không đổi, chủ yếu vì — phát biểu đúng là: thể tích khí tăng theo nhiệt độ.
\item (Sai) Quả bóng mềm căng hơn khi để nơi ấm, nếu áp suất gần không đổi, chủ yếu vì — phát biểu khối lượng khí tăng là đúng.
\end{itemchoice}
}
\end{ex}

% ===== Câu 122 =====
% ID: L12C2B10-31-TL
% Mức: TH
\begin{ex}
% ID: L12C2B10-31-TL
% Mức: TH
Trong dạng “Ứng dụng định luật Charles trong thực tế”, Trình bày cơ sở lí thuyết của dạng “Ứng dụng định luật Charles trong thực tế” và nêu điều kiện áp dụng.
\choiceTF

\loigiai{
Nêu đúng đại lượng không đổi, hệ thức đặc trưng và yêu cầu dùng nhiệt độ kelvin, áp suất tuyệt đối khi có liên quan.
}
\end{ex}

% ===== Câu 123 =====
% ID: L12C2B10-31-TLN
% Mức: VD
\dangbt{Đồ thị đẳng áp trên các hệ tọa độ}
\begin{ex}
% ID: L12C2B10-31-TLN
% Mức: VD
Trong dạng “Đồ thị đẳng áp trên các hệ tọa độ”, Nhiệt độ tuyệt đối tăng 3 lần ở áp suất không đổi. Khối lượng riêng giảm bao nhiêu lần?
\shortans{3}
\loigiai{
$\rho\sim1/T$, nên khối lượng riêng giảm 3 lần.
}
\end{ex}

% ===== Câu 124 =====
% ID: L12C2B10-31-TN
% Mức: VD
\dangbt{Nhận biết quá trình đẳng áp và định luật Charles}
\begin{ex}
% ID: L12C2B10-31-TN
% Mức: VD
Trong dạng “Nhận biết quá trình đẳng áp và định luật Charles”, Trong định luật Charles, nhiệt độ phải đo theo
\choice
{\True kelvin}
{độ $C$ tùy ý}
{độ $F$ tùy ý}
{mọi thang đều được}
\loigiai{
Đáp án A. Tỉ lệ thuận chỉ đúng trực tiếp với nhiệt độ tuyệt đối.
}
\end{ex}

% ===== Câu 125 =====
% ID: L12C2B10-32-DS
% Mức: VDC
\dangbt{Ứng dụng định luật Charles trong thực tế}
\begin{ex}
% ID: L12C2B10-32-DS
% Mức: VDC
Xét các phát biểu sau trong dạng “Ứng dụng định luật Charles trong thực tế”.
\choiceTF
{\True Đường đẳng áp có áp suất nhỏ hơn trên đồ thị V- $T$ của cùng lượng khí nằm — phát biểu đúng là: dốc hơn.}
{Đường đẳng áp có áp suất nhỏ hơn trên đồ thị V- $T$ của cùng lượng khí nằm — phát biểu thoải hơn là đúng.}
{\True Ở áp suất không đổi, thể tích ngoại suy bằng không tại — phát biểu đúng là: 0 K.}
{Ở áp suất không đổi, thể tích ngoại suy bằng không tại — phát biểu 100 ° $C$ là đúng.}
\loigiai{
\begin{itemchoice}
\item (Đúng) Đường đẳng áp có áp suất nhỏ hơn trên đồ thị V- $T$ của cùng lượng khí nằm — phát biểu đúng là: dốc hơn. (Vì): Từ $V=nRT/p$, p nhỏ hơn cho hệ số góc lớn hơn. b) Sai: trái với hệ thức hoặc điều kiện đã nêu. c) Đúng: 0 $K$ là độ không tuyệt đối theo mô hình khí lí tưởng. d) Sai: trái với hệ thức hoặc điều kiện đã nêu
\item (Sai) Đường đẳng áp có áp suất nhỏ hơn trên đồ thị V- $T$ của cùng lượng khí nằm — phát biểu thoải hơn là đúng.
\item (Đúng) Ở áp suất không đổi, thể tích ngoại suy bằng không tại — phát biểu đúng là: 0 K.
\item (Sai) Ở áp suất không đổi, thể tích ngoại suy bằng không tại — phát biểu 100 ° $C$ là đúng.
\end{itemchoice}
}
\end{ex}

% ===== Câu 126 =====
% ID: L12C2B10-32-TL
% Mức: TH
\begin{ex}
% ID: L12C2B10-32-TL
% Mức: TH
Trong dạng “Ứng dụng định luật Charles trong thực tế”, Mô tả dạng đồ thị đặc trưng và cách nhận biết quá trình trong dạng “Ứng dụng định luật Charles trong thực tế”.
\choiceTF

\loigiai{
Xác định các trục, đại lượng không đổi và suy ra hình dạng đồ thị từ hệ thức vật lí tương ứng.
}
\end{ex}

% ===== Câu 127 =====
% ID: L12C2B10-32-TLN
% Mức: VDC
\dangbt{Đồ thị đẳng áp trên các hệ tọa độ}
\begin{ex}
% ID: L12C2B10-32-TLN
% Mức: VDC
Trong dạng “Đồ thị đẳng áp trên các hệ tọa độ”, Khí 16 lít ở 380 $K$ được làm lạnh đẳng áp về 320 K. Tính thể tích theo lít.
\shortans{13,47}
\loigiai{
$V_2=16\cdot320/380=13,47$ lít.
}
\end{ex}

% ===== Câu 128 =====
% ID: L12C2B10-32-TN
% Mức: VD
\dangbt{Nhận biết quá trình đẳng áp và định luật Charles}
\begin{ex}
% ID: L12C2B10-32-TN
% Mức: VD
Trong dạng “Nhận biết quá trình đẳng áp và định luật Charles”, Quả bóng mềm căng hơn khi để nơi ấm, nếu áp suất gần không đổi, chủ yếu vì
\choice
{\True thể tích khí tăng theo nhiệt độ}
{số mol khí tăng}
{khối lượng khí tăng}
{áp suất bằng không}
\loigiai{
Đáp án A. Theo Charles, khi $T$ tăng thì $V$ tăng ở áp suất không đổi.
}
\end{ex}

% ===== Câu 129 =====
% ID: L12C2B10-33-TL
% Mức: VD
\dangbt{Ứng dụng định luật Charles trong thực tế}
\begin{bt}
% ID: L12C2B10-33-TL
% Mức: VD
Trong dạng “Ứng dụng định luật Charles trong thực tế”, Nhiệt độ tuyệt đối tăng từ 310 $K$ lên 370 $K$ ở áp suất không đổi. Tính $V_2/V_1$.
\loigiai{
$V_2/V_1=T_2/T_1=1,19$.
}
\end{bt}

% ===== Câu 130 =====
% ID: L12C2B10-33-TLN
% Mức: TH
\begin{ex}
% ID: L12C2B10-33-TLN
% Mức: TH
Trong dạng “Ứng dụng định luật Charles trong thực tế”, Khí 8 lít ở 290 $K$ được đun đẳng áp đến 350 K. Tính thể tích sau theo lít.
\shortans{9,66}
\loigiai{
$V_2=8\cdot350/290=9,66$ lít.
}
\end{ex}

% ===== Câu 131 =====
% ID: L12C2B10-33-TN
% Mức: VDC
\dangbt{Nhận biết quá trình đẳng áp và định luật Charles}
\begin{ex}
% ID: L12C2B10-33-TN
% Mức: VDC
Trong dạng “Nhận biết quá trình đẳng áp và định luật Charles”, Đường đẳng áp có áp suất nhỏ hơn trên đồ thị V- $T$ của cùng lượng khí nằm
\choice
{\True dốc hơn}
{thoải hơn}
{trùng trục $T$}
{luôn trùng nhau}
\loigiai{
Đáp án A. Từ $V=nRT/p$, p nhỏ hơn cho hệ số góc lớn hơn.
}
\end{ex}

% ===== Câu 132 =====
% ID: L12C2B10-34-TL
% Mức: VD
\dangbt{Ứng dụng định luật Charles trong thực tế}
\begin{ex}
% ID: L12C2B10-34-TL
% Mức: VD
Trong dạng “Ứng dụng định luật Charles trong thực tế”, 4
\choiceTF

\loigiai{
Do $V/T$ không đổi nên $T_2/T_1=V_2/V_1=4$.
}
\end{ex}

% ===== Câu 133 =====
% ID: L12C2B10-34-TLN
% Mức: TH
\begin{ex}
% ID: L12C2B10-34-TLN
% Mức: TH
Trong dạng “Ứng dụng định luật Charles trong thực tế”, Khí có thể tích 4 lít ở 300 $K$, đẳng áp tăng thể tích lên 8 lít. Tính nhiệt độ sau theo K.
\shortans{600}
\loigiai{
$T_2=T_1V_2/V_1=300\cdot2=600$ K.
}
\end{ex}

% ===== Câu 134 =====
% ID: L12C2B10-34-TN
% Mức: VDC
\dangbt{Nhận biết quá trình đẳng áp và định luật Charles}
\begin{ex}
% ID: L12C2B10-34-TN
% Mức: VDC
Trong dạng “Nhận biết quá trình đẳng áp và định luật Charles”, Ở áp suất không đổi, thể tích ngoại suy bằng không tại
\choice
{\True 0 $K$}
{0 ° $C$}
{100 ° $C$}
{273 $K$}
\loigiai{
Đáp án A. 0 $K$ là độ không tuyệt đối theo mô hình khí lí tưởng.
}
\end{ex}

% ===== Câu 135 =====
% ID: L12C2B10-35-TL
% Mức: VD
\dangbt{Ứng dụng định luật Charles trong thực tế}
\begin{bt}
% ID: L12C2B10-35-TL
% Mức: VD
Trong dạng “Ứng dụng định luật Charles trong thực tế”, $\rho\sim1/T$, nên khối lượng riêng giảm 2 lần.
\loigiai{
$\rho\sim1/T$, nên khối lượng riêng giảm 2 lần.
}
\end{bt}

% ===== Câu 136 =====
% ID: L12C2B10-35-TLN
% Mức: VD
\begin{ex}
% ID: L12C2B10-35-TLN
% Mức: VD
Trong dạng “Ứng dụng định luật Charles trong thực tế”, Nhiệt độ tuyệt đối tăng từ 310 $K$ lên 370 $K$ ở áp suất không đổi. Tính $V_2/V_1$.
\shortans{1,19}
\loigiai{
$V_2/V_1=T_2/T_1=1,19$.
}
\end{ex}

% ===== Câu 137 =====
% ID: L12C2B10-35-TN
% Mức: NB
\dangbt{Tính thể tích và nhiệt độ trong quá trình đẳng áp}
\begin{ex}
% ID: L12C2B10-35-TN
% Mức: NB
Trong dạng “Tính thể tích và nhiệt độ trong quá trình đẳng áp”, Trong quá trình đẳng áp của một lượng khí xác định, đại lượng không đổi là
\choice
{\True áp suất}
{thể tích}
{nhiệt độ}
{khối lượng riêng}
\loigiai{
Đáp án A. Quá trình đẳng áp có áp suất không đổi.
}
\end{ex}

% ===== Câu 138 =====
% ID: L12C2B10-36-TL
% Mức: VDC
\dangbt{Ứng dụng định luật Charles trong thực tế}
\begin{bt}
% ID: L12C2B10-36-TL
% Mức: VDC
Trong dạng “Ứng dụng định luật Charles trong thực tế”, undefined
\loigiai{
$V_2=18\cdot290/350=14,91$ lít.
}
\end{bt}

% ===== Câu 139 =====
% ID: L12C2B10-36-TLN
% Mức: VD
\begin{ex}
% ID: L12C2B10-36-TLN
% Mức: VD
Trong dạng “Ứng dụng định luật Charles trong thực tế”, Thể tích đẳng áp tăng 4 lần. Nhiệt độ tuyệt đối tăng bao nhiêu lần?
\shortans{4}
\loigiai{
Do $V/T$ không đổi nên $T_2/T_1=V_2/V_1=4$.
}
\end{ex}

% ===== Câu 140 =====
% ID: L12C2B10-36-TN
% Mức: NB
\dangbt{Tính thể tích và nhiệt độ trong quá trình đẳng áp}
\begin{ex}
% ID: L12C2B10-36-TN
% Mức: NB
Trong dạng “Tính thể tích và nhiệt độ trong quá trình đẳng áp”, Hệ thức đúng của định luật Charles là
\choice
{\True $V_1/T_1=V_2/T_2$}
{$p_1V_1=p_2V_2$}
{$p_1/T_1=p_2/T_2$}
{$V_1T_1=V_2T_2$}
\loigiai{
Đáp án A. Ở áp suất không đổi, $V/T$ là hằng số.
}
\end{ex}

% ===== Câu 141 =====
% ID: L12C2B10-37-TLN
% Mức: VD
\dangbt{Ứng dụng định luật Charles trong thực tế}
\begin{ex}
% ID: L12C2B10-37-TLN
% Mức: VD
Trong dạng “Ứng dụng định luật Charles trong thực tế”, Nhiệt độ tuyệt đối tăng 2 lần ở áp suất không đổi. Khối lượng riêng giảm bao nhiêu lần?
\shortans{2}
\loigiai{
$\rho\sim1/T$, nên khối lượng riêng giảm 2 lần.
}
\end{ex}

% ===== Câu 142 =====
% ID: L12C2B10-37-TN
% Mức: TH
\dangbt{Tính thể tích và nhiệt độ trong quá trình đẳng áp}
\begin{ex}
% ID: L12C2B10-37-TN
% Mức: TH
Trong dạng “Tính thể tích và nhiệt độ trong quá trình đẳng áp”, Khí có thể tích 7 lít ở 270 $K$, đun đẳng áp đến 340 K. Thể tích sau gần bằng
\choice
{\True 8,81 lít}
{5,56 lít}
{7 lít}
{17,63 lít}
\loigiai{
Đáp án A. $V_2=V_1T_2/T_1=7\cdot340/270=8,81$ lít.
}
\end{ex}

% ===== Câu 143 =====
% ID: L12C2B10-38-TLN
% Mức: VDC
\dangbt{Ứng dụng định luật Charles trong thực tế}
\begin{ex}
% ID: L12C2B10-38-TLN
% Mức: VDC
Trong dạng “Ứng dụng định luật Charles trong thực tế”, Khí 18 lít ở 350 $K$ được làm lạnh đẳng áp về 290 K. Tính thể tích theo lít.
\shortans{14,91}
\loigiai{
$V_2=18\cdot290/350=14,91$ lít.
}
\end{ex}

% ===== Câu 144 =====
% ID: L12C2B10-38-TN
% Mức: TH
\dangbt{Tính thể tích và nhiệt độ trong quá trình đẳng áp}
\begin{ex}
% ID: L12C2B10-38-TN
% Mức: TH
Trong dạng “Tính thể tích và nhiệt độ trong quá trình đẳng áp”, Khí đẳng áp tăng nhiệt độ tuyệt đối từ 280 $K$ lên 360 K. Tỉ số $V_2/V_1$ bằng
\choice
{\True 1,29}
{0,78}
{1}
{0,29}
\loigiai{
Đáp án A. $V_2/V_1=T_2/T_1=1,29$.
}
\end{ex}

% ===== Câu 145 =====
% ID: L12C2B10-39-TN
% Mức: TH
\begin{ex}
% ID: L12C2B10-39-TN
% Mức: TH
Trong dạng “Tính thể tích và nhiệt độ trong quá trình đẳng áp”, Đồ thị V- $T$ của quá trình đẳng áp trong thang kelvin là
\choice
{\True đường thẳng qua gốc}
{hypebol}
{đường ngang}
{đường thẳng không qua gốc}
\loigiai{
Đáp án A. $V$ tỉ lệ thuận $T$ tuyệt đối.
}
\end{ex}

% ===== Câu 146 =====
% ID: L12C2B10-40-TN
% Mức: VD
\begin{ex}
% ID: L12C2B10-40-TN
% Mức: VD
Trong dạng “Tính thể tích và nhiệt độ trong quá trình đẳng áp”, Khi áp suất không đổi và nhiệt độ tuyệt đối tăng gấp đôi, khối lượng riêng
\choice
{\True giảm một nửa}
{tăng gấp đôi}
{không đổi}
{tăng bốn lần}
\loigiai{
Đáp án A. $\rho=m/V$ và $V\sim T$, nên $\rho\sim1/T$.
}
\end{ex}

% ===== Câu 147 =====
% ID: L12C2B10-41-TN
% Mức: VD
\begin{ex}
% ID: L12C2B10-41-TN
% Mức: VD
Trong dạng “Tính thể tích và nhiệt độ trong quá trình đẳng áp”, Trong định luật Charles, nhiệt độ phải đo theo
\choice
{\True kelvin}
{độ $C$ tùy ý}
{độ $F$ tùy ý}
{mọi thang đều được}
\loigiai{
Đáp án A. Tỉ lệ thuận chỉ đúng trực tiếp với nhiệt độ tuyệt đối.
}
\end{ex}

% ===== Câu 148 =====
% ID: L12C2B10-42-TN
% Mức: VD
\begin{ex}
% ID: L12C2B10-42-TN
% Mức: VD
Trong dạng “Tính thể tích và nhiệt độ trong quá trình đẳng áp”, Quả bóng mềm căng hơn khi để nơi ấm, nếu áp suất gần không đổi, chủ yếu vì
\choice
{\True thể tích khí tăng theo nhiệt độ}
{số mol khí tăng}
{khối lượng khí tăng}
{áp suất bằng không}
\loigiai{
Đáp án A. Theo Charles, khi $T$ tăng thì $V$ tăng ở áp suất không đổi.
}
\end{ex}

% ===== Câu 149 =====
% ID: L12C2B10-43-TN
% Mức: VDC
\begin{ex}
% ID: L12C2B10-43-TN
% Mức: VDC
Trong dạng “Tính thể tích và nhiệt độ trong quá trình đẳng áp”, Đường đẳng áp có áp suất nhỏ hơn trên đồ thị V- $T$ của cùng lượng khí nằm
\choice
{\True dốc hơn}
{thoải hơn}
{trùng trục $T$}
{luôn trùng nhau}
\loigiai{
Đáp án A. Từ $V=nRT/p$, p nhỏ hơn cho hệ số góc lớn hơn.
}
\end{ex}

% ===== Câu 150 =====
% ID: L12C2B10-44-TN
% Mức: VDC
\begin{ex}
% ID: L12C2B10-44-TN
% Mức: VDC
Trong dạng “Tính thể tích và nhiệt độ trong quá trình đẳng áp”, Ở áp suất không đổi, thể tích ngoại suy bằng không tại
\choice
{\True 0 $K$}
{0 ° $C$}
{100 ° $C$}
{273 $K$}
\loigiai{
Đáp án A. 0 $K$ là độ không tuyệt đối theo mô hình khí lí tưởng.
}
\end{ex}

% ===== Câu 151 =====
% ID: L12C2B10-45-TN
% Mức: NB
\dangbt{Đồ thị đẳng áp trên các hệ tọa độ}
\begin{ex}
% ID: L12C2B10-45-TN
% Mức: NB
Trong dạng “Đồ thị đẳng áp trên các hệ tọa độ”, Trong quá trình đẳng áp của một lượng khí xác định, đại lượng không đổi là
\choice
{\True áp suất}
{thể tích}
{nhiệt độ}
{khối lượng riêng}
\loigiai{
Đáp án A. Quá trình đẳng áp có áp suất không đổi.
}
\end{ex}

% ===== Câu 152 =====
% ID: L12C2B10-46-TN
% Mức: NB
\begin{ex}
% ID: L12C2B10-46-TN
% Mức: NB
Trong dạng “Đồ thị đẳng áp trên các hệ tọa độ”, Hệ thức đúng của định luật Charles là
\choice
{\True $V_1/T_1=V_2/T_2$}
{$p_1V_1=p_2V_2$}
{$p_1/T_1=p_2/T_2$}
{$V_1T_1=V_2T_2$}
\loigiai{
Đáp án A. Ở áp suất không đổi, $V/T$ là hằng số.
}
\end{ex}

% ===== Câu 153 =====
% ID: L12C2B10-47-TN
% Mức: TH
\begin{ex}
% ID: L12C2B10-47-TN
% Mức: TH
Trong dạng “Đồ thị đẳng áp trên các hệ tọa độ”, Khí có thể tích 8 lít ở 280 $K$, đun đẳng áp đến 350 K. Thể tích sau gần bằng
\choice
{\True 10 lít}
{6,4 lít}
{8 lít}
{20 lít}
\loigiai{
Đáp án A. $V_2=V_1T_2/T_1=8\cdot350/280=10$ lít.
}
\end{ex}

% ===== Câu 154 =====
% ID: L12C2B10-48-TN
% Mức: TH
\begin{ex}
% ID: L12C2B10-48-TN
% Mức: TH
Trong dạng “Đồ thị đẳng áp trên các hệ tọa độ”, Khí đẳng áp tăng nhiệt độ tuyệt đối từ 290 $K$ lên 370 K. Tỉ số $V_2/V_1$ bằng
\choice
{\True 1,28}
{0,78}
{1}
{0,28}
\loigiai{
Đáp án A. $V_2/V_1=T_2/T_1=1,28$.
}
\end{ex}

% ===== Câu 155 =====
% ID: L12C2B10-49-TN
% Mức: TH
\begin{ex}
% ID: L12C2B10-49-TN
% Mức: TH
Trong dạng “Đồ thị đẳng áp trên các hệ tọa độ”, Đồ thị V- $T$ của quá trình đẳng áp trong thang kelvin là
\choice
{\True đường thẳng qua gốc}
{hypebol}
{đường ngang}
{đường thẳng không qua gốc}
\loigiai{
Đáp án A. $V$ tỉ lệ thuận $T$ tuyệt đối.
}
\end{ex}

% ===== Câu 156 =====
% ID: L12C2B10-50-TN
% Mức: VD
\begin{ex}
% ID: L12C2B10-50-TN
% Mức: VD
Trong dạng “Đồ thị đẳng áp trên các hệ tọa độ”, Khi áp suất không đổi và nhiệt độ tuyệt đối tăng gấp đôi, khối lượng riêng
\choice
{\True giảm một nửa}
{tăng gấp đôi}
{không đổi}
{tăng bốn lần}
\loigiai{
Đáp án A. $\rho=m/V$ và $V\sim T$, nên $\rho\sim1/T$.
}
\end{ex}

% ===== Câu 157 =====
% ID: L12C2B10-51-TN
% Mức: VD
\begin{ex}
% ID: L12C2B10-51-TN
% Mức: VD
Trong dạng “Đồ thị đẳng áp trên các hệ tọa độ”, Trong định luật Charles, nhiệt độ phải đo theo
\choice
{\True kelvin}
{độ $C$ tùy ý}
{độ $F$ tùy ý}
{mọi thang đều được}
\loigiai{
Đáp án A. Tỉ lệ thuận chỉ đúng trực tiếp với nhiệt độ tuyệt đối.
}
\end{ex}

% ===== Câu 158 =====
% ID: L12C2B10-52-TN
% Mức: VD
\begin{ex}
% ID: L12C2B10-52-TN
% Mức: VD
Trong dạng “Đồ thị đẳng áp trên các hệ tọa độ”, Quả bóng mềm căng hơn khi để nơi ấm, nếu áp suất gần không đổi, chủ yếu vì
\choice
{\True thể tích khí tăng theo nhiệt độ}
{số mol khí tăng}
{khối lượng khí tăng}
{áp suất bằng không}
\loigiai{
Đáp án A. Theo Charles, khi $T$ tăng thì $V$ tăng ở áp suất không đổi.
}
\end{ex}

% ===== Câu 159 =====
% ID: L12C2B10-53-TN
% Mức: VDC
\begin{ex}
% ID: L12C2B10-53-TN
% Mức: VDC
Trong dạng “Đồ thị đẳng áp trên các hệ tọa độ”, Đường đẳng áp có áp suất nhỏ hơn trên đồ thị V- $T$ của cùng lượng khí nằm
\choice
{\True dốc hơn}
{thoải hơn}
{trùng trục $T$}
{luôn trùng nhau}
\loigiai{
Đáp án A. Từ $V=nRT/p$, p nhỏ hơn cho hệ số góc lớn hơn.
}
\end{ex}

% ===== Câu 160 =====
% ID: L12C2B10-54-TN
% Mức: VDC
\begin{ex}
% ID: L12C2B10-54-TN
% Mức: VDC
Trong dạng “Đồ thị đẳng áp trên các hệ tọa độ”, Ở áp suất không đổi, thể tích ngoại suy bằng không tại
\choice
{\True 0 $K$}
{0 ° $C$}
{100 ° $C$}
{273 $K$}
\loigiai{
Đáp án A. 0 $K$ là độ không tuyệt đối theo mô hình khí lí tưởng.
}
\end{ex}

% ===== Câu 161 =====
% ID: L12C2B10-55-TN
% Mức: NB
\dangbt{Ứng dụng định luật Charles trong thực tế}
\begin{ex}
% ID: L12C2B10-55-TN
% Mức: NB
Trong dạng “Ứng dụng định luật Charles trong thực tế”, Trong quá trình đẳng áp của một lượng khí xác định, đại lượng không đổi là
\choice
{\True áp suất}
{thể tích}
{nhiệt độ}
{khối lượng riêng}
\loigiai{
Đáp án A. Quá trình đẳng áp có áp suất không đổi.
}
\end{ex}

% ===== Câu 162 =====
% ID: L12C2B10-56-TN
% Mức: NB
\begin{ex}
% ID: L12C2B10-56-TN
% Mức: NB
Trong dạng “Ứng dụng định luật Charles trong thực tế”, Hệ thức đúng của định luật Charles là
\choice
{\True $V_1/T_1=V_2/T_2$}
{$p_1V_1=p_2V_2$}
{$p_1/T_1=p_2/T_2$}
{$V_1T_1=V_2T_2$}
\loigiai{
Đáp án A. Ở áp suất không đổi, $V/T$ là hằng số.
}
\end{ex}

% ===== Câu 163 =====
% ID: L12C2B10-57-TN
% Mức: TH
\begin{ex}
% ID: L12C2B10-57-TN
% Mức: TH
Trong dạng “Ứng dụng định luật Charles trong thực tế”, Khí có thể tích 3 lít ở 300 $K$, đun đẳng áp đến 370 K. Thể tích sau gần bằng
\choice
{\True 3,7 lít}
{2,43 lít}
{3 lít}
{7,4 lít}
\loigiai{
Đáp án A. $V_2=V_1T_2/T_1=3\cdot370/300=3,7$ lít.
}
\end{ex}

% ===== Câu 164 =====
% ID: L12C2B10-58-TN
% Mức: TH
\begin{ex}
% ID: L12C2B10-58-TN
% Mức: TH
Trong dạng “Ứng dụng định luật Charles trong thực tế”, Khí đẳng áp tăng nhiệt độ tuyệt đối từ 310 $K$ lên 390 K. Tỉ số $V_2/V_1$ bằng
\choice
{\True 1,26}
{0,79}
{1}
{0,26}
\loigiai{
Đáp án A. $V_2/V_1=T_2/T_1=1,26$.
}
\end{ex}

% ===== Câu 165 =====
% ID: L12C2B10-59-TN
% Mức: TH
\begin{ex}
% ID: L12C2B10-59-TN
% Mức: TH
Trong dạng “Ứng dụng định luật Charles trong thực tế”, Đồ thị V- $T$ của quá trình đẳng áp trong thang kelvin là
\choice
{\True đường thẳng qua gốc}
{hypebol}
{đường ngang}
{đường thẳng không qua gốc}
\loigiai{
Đáp án A. $V$ tỉ lệ thuận $T$ tuyệt đối.
}
\end{ex}

% ===== Câu 166 =====
% ID: L12C2B10-60-TN
% Mức: VD
\begin{ex}
% ID: L12C2B10-60-TN
% Mức: VD
Trong dạng “Ứng dụng định luật Charles trong thực tế”, Khi áp suất không đổi và nhiệt độ tuyệt đối tăng gấp đôi, khối lượng riêng
\choice
{\True giảm một nửa}
{tăng gấp đôi}
{không đổi}
{tăng bốn lần}
\loigiai{
Đáp án A. $\rho=m/V$ và $V\sim T$, nên $\rho\sim1/T$.
}
\end{ex}

% ===== Câu 167 =====
% ID: L12C2B10-61-TN
% Mức: VD
\begin{ex}
% ID: L12C2B10-61-TN
% Mức: VD
Trong dạng “Ứng dụng định luật Charles trong thực tế”, Trong định luật Charles, nhiệt độ phải đo theo
\choice
{\True kelvin}
{độ $C$ tùy ý}
{độ $F$ tùy ý}
{mọi thang đều được}
\loigiai{
Đáp án A. Tỉ lệ thuận chỉ đúng trực tiếp với nhiệt độ tuyệt đối.
}
\end{ex}

% ===== Câu 168 =====
% ID: L12C2B10-62-TN
% Mức: VD
\begin{ex}
% ID: L12C2B10-62-TN
% Mức: VD
Trong dạng “Ứng dụng định luật Charles trong thực tế”, Quả bóng mềm căng hơn khi để nơi ấm, nếu áp suất gần không đổi, chủ yếu vì
\choice
{\True thể tích khí tăng theo nhiệt độ}
{số mol khí tăng}
{khối lượng khí tăng}
{áp suất bằng không}
\loigiai{
Đáp án A. Theo Charles, khi $T$ tăng thì $V$ tăng ở áp suất không đổi.
}
\end{ex}

% ===== Câu 169 =====
% ID: L12C2B10-63-TN
% Mức: VDC
\begin{ex}
% ID: L12C2B10-63-TN
% Mức: VDC
Trong dạng “Ứng dụng định luật Charles trong thực tế”, Đường đẳng áp có áp suất nhỏ hơn trên đồ thị V- $T$ của cùng lượng khí nằm
\choice
{\True dốc hơn}
{thoải hơn}
{trùng trục $T$}
{luôn trùng nhau}
\loigiai{
Đáp án A. Từ $V=nRT/p$, p nhỏ hơn cho hệ số góc lớn hơn.
}
\end{ex}

% ===== Câu 170 =====
% ID: L12C2B10-64-TN
% Mức: VDC
\begin{ex}
% ID: L12C2B10-64-TN
% Mức: VDC
Trong dạng “Ứng dụng định luật Charles trong thực tế”, Ở áp suất không đổi, thể tích ngoại suy bằng không tại
\choice
{\True 0 $K$}
{0 ° $C$}
{100 ° $C$}
{273 $K$}
\loigiai{
Đáp án A. 0 $K$ là độ không tuyệt đối theo mô hình khí lí tưởng.
}
\end{ex}
